\documentclass[a4paper, english, 12pt, bibliography=totocnumbered,
index=totoc]{scrartcl}

\usepackage{mystyles}

\title{The Chern character of \texorpdfstring{$\vartheta$}{ϑ}-summable
\texorpdfstring{$\mathcal{C}_q$}{C(q)}-Fredholm modules over locally convex
differential graded algebras}
\author{Jonas Miehe}
\date{December 11, 2023}

\begin{document}
\maketitle

\begin{abstract}
	Let $\Omega$ be a locally convex \emph{differential graded} algebra.
	We introduce the Chern character of $\vartheta$-summable
	$\mathcal{C}_q$-Fredholm modules over $\Omega$,
	generalizing the JLO cocycle to the differential graded setting.
	This allows treating the case of even $\vartheta$-summable Fredholm modules
	over $\Omega$
	(which has been previously considered by Güneysu/Ludewig in the context of
	localization on loop spaces) and odd
	$\vartheta$-summable Fredholm modules over $\Omega$
	simultaneously.
	We prove all naturally expected properties of this construction and
	relate it to the spectral flow in the odd case.
\end{abstract}
%\tableofcontents

\section{Introduction}
Over the last 35 years, deep connections between index theory, noncommutative
geometry and cyclic homology have been shown to exist using the
algebraic framework provided by Fredholm modules.
Given an even $\vartheta$-summable Fredholm module over a unital Banach
$*$-algebra, it was shown by Connes, \cite{ConnesCyclicCohomology}, and
Jaffe, Lesniewski and Osterwalder, \cite{JLO}, that there exists a Chern
character, an entire cyclic
cocycle on $\mathcal{A}$, which admits a pairing with $K$-theory
$K_0(\mathcal{A})$ thus producing a noncommutative index theorem,
\cite{GetzlerSzenes}, generalizing previous results obtained in the
finitely-summable case.
The weaker $\vartheta$-summability assumption and the
fact that one can allow for infinite chains (subject to entire growth
conditions) makes it possible to treat infinite-dimensional situations, such as
applications in quantum field theory.

This program was completed by Getzler in \cite{GetzlerOdd} who introduced the
Chern character of an \emph{odd} $\vartheta$-summable Fredholm module --
\textit{the odd JLO cocycle} -- which can be
paired with odd $K$-theory $K_1(\mathcal{A})$ producing a natural replacement
of the noncommutative index theorem in the odd case, where the index is
replaced with the spectral flow.

Motivated by an infinite-dimensional variant of Duistermaat-Heckman
localization in equivariant loop space cohomology, which allows for an
remarkably short and elegant -- however formal --
proof of the Atiyah-Singer index theorem, Güneysu and Ludewig considerably
extended the theory of even $\vartheta$-summable Fredholm modules to the
setting of locally convex \emph{differential graded} algebras,
\cite{GüneysuLudewig}.
They constructed an analytic cyclic cocycle which admits a
noncommutative index theorem, which is invariant under suitable homotopies of
Fredholm modules, and which coincides with the JLO cocycle in case the algebra
is trivially graded with the trivial differential.
These constructions allow formulating and proving a localization formula on the
loop space via Chen's iterated integrals in mathematically rigorous way,
confirming the conjectures made by Atiyah, Bismut and Witten
\cite{AtiyahCircular,BismutCMP85}.

The goal of this paper is to complete this program, which can be traced back at
least to Getzler/Jones/Petrack, \cite{GetzlerJonesPetrack}, by extending these
results to the odd-dimensional case.
To this end, given $q\in\mathbb{N}_0$, we will introduce the notion of a
\emph{$\vartheta$-summable $\mathcal{C}_q$-Fredholm module $\mathscr{M}$} over
a locally convex differential graded algebra $\Omega$, where an additional
action of the Clifford algebra $\mathcal{C}_q$ in fact allows us to treat the
even and the (so far not at all treated) odd case
simultaneously: every even $\vartheta$-summable Fredholm module over $\Omega$
is an $\mathcal{C}_0$-Fredholm module over $\Omega$, and every odd
$\vartheta$-summable Fredholm module $\mathscr{M}$ over $\Omega$ canonically
induces an $\vartheta$-summable $\mathcal{C}_1$-Fredholm module
$\tilde{\mathscr{M}}$ over $\Omega$.
Our first main result reads as follows:

\renewcommand*{\thethmALPH}{\Alph{thmALPH}}
\begin{thmALPH}
	Let $\mathscr{M}$ be a $\vartheta$-summable $\mathcal{C}_q$-Fredholm module
	over a locally convex differential graded algebra $\Omega$.
	There exists a canonically given analytic cochain on the entire cyclic
	complex over $\Omega$, the \emph{Chern character $\Chern$ of
	$\mathscr{M}$}, with the following properties:
	\begin{enumerate}
		\item $\Chern$ is coclosed.
		\item If $\mathscr{M}_\T$ denotes the acyclic extension of
		$\mathscr{M}$ (a module over the acyclic extension $\Omega_\T$ of
		$\Omega$), then $\ChernT$ vanishes on the Chen normalization of the
		entire cyclic complex over $\Omega_\T$.
	\end{enumerate}
	If $\mathscr{M}_s$, $s \in [0,1]$, is a homotopy of $\vartheta$-summable
	$\mathcal{C}_q$-Fredholm modules over $\Omega$, then
	$\mathrm{Ch}_{\mathscr{M}^0_\T}$ and
	$\mathrm{Ch}_{\mathscr{M}^1_\T}$ are homologous in the Chen normalization
	of the entire cyclic complex over $\Omega_\T$.
\end{thmALPH}
Also, we show that if $\mathscr{M}$ is an odd $\vartheta$-summable weak
$\mathcal{C}_q$-Fredholm module over a locally convex differential graded
algebra $\Omega$, then the Chern character of the induced
$\mathcal{C}_1$-Fredholm module
$\tilde{\mathscr{M}}$ over $\Omega$ is equal to the odd JLO cocycle in case
$\Omega$ is trivially graded with trivial differential.

In \cref{Section: definition cherng} we generalize the constructions from
\cite{GüneysuCacciatori} to arbitrary
locally convex differential graded algebras thereby proving:
\begin{thmALPH}
	Assume $\Omega$ is a locally convex differential graded algebra.
	For every $g \in \mathrm{GL}_m(\Omega^0)$ there exists a canonically given
	entire cyclic chain $\mathrm{Ch}(g)$ of the entire cyclic complex over
	$\Omega_\T$; this chain is closed in the Chen normalization of the entire
	cyclic complex over $\Omega_\T$.
\end{thmALPH}
Using the combinatorial machinery of the Chern character, we prove a
noncommutative index theorem by pairing the Chern character with the entire
chain from above:
\begin{thmALPH}
	Let $\mathscr{M}$ be an odd $\vartheta$-summable Fredholm
	module over a locally convex differential graded algebra $\Omega$.
	Then one has
	\begin{equation}\label{formel1}
		\left\langle \mathrm{Ch}_{\widetilde{\mathscr{M}_\T}
		},\mathrm{Ch}(g)\right\rangle =\int_0^1 \Tr(\dot{Q}_s e^{-Q_s^2}) \,
		\mathrm{d}s
	\end{equation}
	where $Q_s = (1-s) Q + s \mathbf{c}(g^{-1}) Q \mathbf{c}(g)$ and
	$\dot{Q}_s= \frac{d}{ds}Q_s$.
\end{thmALPH}
The left hand side of \cref{formel1} is precisely the spectral flow
$\mathrm{sf}(Q, \mathbf{c}(g^{-1}) Q
\mathbf{c}(g))$ if $\mathbf{c}(g)$ is unitary.

This paper is organized as follows: in \cref{Section: Algebraic
preliminaries} we give an overview of the algebraic constructions connected
to the cyclic complex and the bar complex.
Moreover, we recall the notion of analytic chains and entire cochains and
define the Chen normalized complex.
\cref{Section: Cq FM} will be devoted to defining $\vartheta$-summable
$\mathcal{C}_q$-Fredholm modules and homotopies of them.
After that, in \cref{Section: odd FM}, we will see how odd
$\vartheta$-summable Fredholm modules fit into this more general framework.
In \cref{Section: Chern character} we will construct the Chern character of a
$\vartheta$-summable $\mathcal{C}_q$-Fredholm module as a cyclic cocycle and
present its properties.
Moreover, we will see that in the odd case this is indeed a differential
graded extension of the odd JLO cocycle.
Finally, in \cref{Section: definition cherng}, we construct the chains
$\mathrm{Ch}(g)$ and prove \cref{formel1} for the pairing with
the Chern character.

\textbf{Acknowledgements:} This paper is part of the author's PhD thesis under
the supervision of Prof. \!Dr. \!Batu Güneysu at the Technical University of
Chemnitz, Germany.
The author would like to thank Sergio Cacciatori and Shu Shen for several
helpful discussions.
\section{Algebraic preliminaries}\label{Section: Algebraic preliminaries}
All vector spaces and algebras will be assumed to be over $\C$, and we always
work with \textit{graded} vector spaces unless specified otherwise.
On a $\Z$-graded vector space we will use the canonical $\Z_2$-grading 
determined by the even and odd elements and correspondingly, tensor products 
and direct sums of $\Z$-graded vector spaces are canonically $\Z$-graded and 
thereby $\Z_2$-graded.
If a vector space carries no specified grading, we turn it into a 
$\Z$-graded (and therefore $\Z_2$-graded) vector space by letting all 
elements have degree zero.

Given $\Z_2$-graded vector spaces $V$ and $W$, the vector space
$\Hom(V, W)$ is $\Z_2$-graded as usual.
Any $A \in \End(V)$ gives rise to a $\Z_2$-graded dual map $A^{\vee} \colon 
\Hom(V, W) \to \Hom(V, W)$ with values in $W$ given by
\begin{align*}
	(A^{\vee}\ell) (v) \coloneqq (-1)^{|A||\ell|}\ell (A(v)), \quad \ell\in
	\mathrm{Hom}(V,W), \quad v\in V.
\end{align*}
We use this definition of the dual map instead of the usual one unless 
specified otherwise.
For $A, B\in \End(V)$ we will denote by $[A, B] \in \End(V)$
the graded commutator (or supercommutator), i. e.
\begin{equation*}
	[A,B] \coloneqq AB-(-1)^{|A||B|}BA,
\end{equation*}
and we will use the notation $[[A, B]] \coloneqq AB + BA$ for the 
anticommutator.
\subsection{Differential graded algebras}
\begin{defn}
	A \textit{differential graded algebra} (\enquote{dg algebra}) $\Omega$ is
	an associative, $\Z$-graded algebra with unit $\mathbf{1}$ equipped
	with a degree $+1$ differential $d$ subject to the graded Leibniz
	rule
	\begin{equation*}
		d(ab) = (da)b + (-1)^{\lvert a \rvert}adb
	\end{equation*}
	for  $a, b \in \Omega$ of pure degree.

	If $\Omega$ carries a locally convex topology\footnote{We always assume
		locally convex topologies to be Hausdorff.} such that $d$ is
	continuous, multiplication is jointly continuous and $\Omega$ is the
	topological direct
	sum of its homogeneous subspaces, then $\Omega$ is a \textit{locally
		convex dg algebra}.
\end{defn}
\begin{eg}
	Let $X$ be a smooth manifold and $\Omega \coloneqq \Omega(X)$ be the
	space of differential forms on $X$.
	Together with the exterior differential and the wedge product, $\Omega$
	is a dg algebra which carries a canonical Fréchet topology turning it
	into a locally convex dg algebra.
\end{eg}
The following construction will be important later on:
\begin{eg}\label{Acyclic extension}
	Let $\Omega$ be a locally convex dg algebra.
	Denote by
	\begin{equation*}
		\Omega_{\mathbb{T}} \coloneqq \Omega[\sigma]
	\end{equation*}
	the $\Z$-graded unital algebra obtained by adjoining the formal variable
	$\sigma$ of degree $-1$ with $\sigma^2 = 0$ and supercommuting with
	elements of $\Omega$.
	By definition, $\theta \in \Omega_{\mathbb{T}}$ can be uniquely written as
	$\theta = \theta' + \sigma \theta''$ where $\theta', \theta'' \in \Omega$.
	We define a differential $d_{\mathbb{T}}$ via
	\begin{equation*}
		d_{\mathbb{T}} \coloneqq d - \iota \quad \text{where} \quad d \theta
		\coloneqq
		d \theta' - \sigma d \theta'' \quad \text{and} \quad \iota \theta
		\coloneqq
		\theta''.
	\end{equation*}
	Using the isomorphism of graded vector spaces $\Omega_{\mathbb{T}} \simeq
	\Omega \oplus \Omega[1]$ (where $\Omega[1]$ is equal to $\Omega$ with
	degrees shifted upwards by one), it follows that $\Omega_{\mathbb{T}}$
	becomes a locally convex dg algebra which we will call the
	\textit{acyclic extension of $\Omega$}.
\end{eg}
Throughout the rest of this chapter, let $\Omega$ be a locally convex dg
algebra.
\subsection{The (reduced) cyclic complex}\label{Section: cyclic complex}
\textit{The (reduced) cyclic complex associated to $\Omega$} is the vector 
space
\begin{equation*}
	\mathrm{C}(\Omega) \coloneqq \bigoplus_{N=0}^{\infty}\Omega \otimes
	\underline{\Omega}[1]^{\otimes N}
\end{equation*}
with grading determined by
\begin{equation*}
	\mathrm{C}_n(\Omega) = \bigoplus_{N=0}^{\infty} \bigoplus_{l_0 + \ldots +
		l_N = n + N} \Omega^{l_0} \otimes \underline{\Omega}^{l_1} \otimes
	\ldots
	\otimes \underline{\Omega}^{l_N}
\end{equation*}
for $n \in \Z$ where $\underline{\Omega} \coloneqq \Omega / \C \mathbf{1}$ is
the quotient vector space.
We will usually denote elements of $\mathrm{C}(\Omega)$ by $(\theta_0,
\ldots, \theta_N)$ instead of $\theta_0 \otimes \ldots \otimes \theta_N$.
We define two degree $+1$ differentials on $\mathrm{C}(\Omega)$ via
\begin{align}
	\underline{d}(\theta_0, \ldots, \theta_N) & \coloneqq (d
	\theta_0, \theta_1, \ldots, \theta_N) - \sum_{k=1}^N
	(-1)^{m_{k-1}}(\theta_0, \ldots, d
	\theta_k, \ldots, \theta_N) \notag \\
	\begin{split}\label{Definition bunderline}
		\underline{b}(\theta_0, \ldots, \theta_N) & \coloneqq
		-\sum_{k=0}^{N-1}(-1)^{m_{k}}(\theta_0, \ldots, \theta_k
		\theta_{k+1},
		\ldots, \theta_N) \\
		& \qquad + (-1)^{(\lvert \theta_{N} \rvert - 1)
			m_{N-1}}(\theta_N\theta_0, \ldots,\theta_{N-1})
	\end{split}
\end{align}
where $m_k \coloneqq \lvert \theta_0 \rvert + \ldots + \lvert \theta_k \rvert
- k$.

The \textit{Connes operator on $\mathrm{C}(\Omega)$} is the degree $-1$
differential defined via
\begin{equation*}
	\underline{B}(\theta_0, \ldots, \theta_N) \coloneqq \sum_{k=0}^N
	(-1)^{(m_{k-1} + 1)(m_N - m_{k-1})}(\mathbf{1}, \theta_{k},
	\ldots,
	\theta_{N}, \theta_0, \ldots, \theta_{k-1})
\end{equation*}
where $m_{-1} = 1$.
The following statement follows from a standard calculation:
\begin{lem}
	The maps $\underline{d}$, $\underline{b}$ and $\underline{B}$ are
	well-defined, pairwise anticommuting differentials.
\end{lem}
Let $\mathrm{C}_{+}(\Omega)$ be the subspace of $\mathrm{C}(\Omega)$ 
containing all elements of even 
degree and $\mathrm{C}_{-}(\Omega)$ be the subspaces containing all elements 
of odd degree.
It follows that we have a $\Z_2$-graded complex
\begin{center}
	\begin{tikzcd}
		\mathrm{C}_{+}(\Omega) \arrow[rr,
		"\underline{d}+\underline{b}-\underline{B}",
		bend left] &  & \mathrm{C}_{-}(\Omega) \arrow[ll,
		"\underline{d}+\underline{b}-\underline{B}", bend left]
	\end{tikzcd}.
\end{center}
\subsection{The bar construction}\label{Section: bar construction}
\subsubsection{Bar chains}
The \textit{bar construction for $\Omega$} is defined as the vector
space
\begin{equation}\label{Bar construction}
	\mathrm{B}(\Omega) \coloneqq \bigoplus_{N=0}^{\infty} \Omega[1]^{\otimes
		N}
\end{equation}
with grading given by
\begin{equation*}
	\mathrm{B}_n(\Omega) = \bigoplus_{N=0}^{\infty} \bigoplus_{l_1 + \ldots +
		l_N = n + N} \Omega^{l_1} \otimes \ldots \otimes \Omega^{l_N}
\end{equation*}
for $n \in \Z$.
We will usually denote elements of $\mathrm{B}(\Omega)$ by $(\theta_1,
\ldots, \theta_N)$ instead of $\theta_1 \otimes \ldots \otimes \theta_N$.
On $\mathrm{B}(\Omega)$ we have two differentials of degree $+1$ given by
\begin{align*}
	d(\theta_1, \ldots, \theta_N) & \coloneqq -\sum_{k=1}^N (-1)^{n_{k-1}}
	(\theta_1, \ldots, d \theta_k, \ldots, \theta_N) \\
	b'(\theta_1, \ldots, \theta_N) & \coloneqq - \sum_{k=1}^{N-1} (-1)^{n_{k}}
	(\theta_1, \ldots, \theta_k \theta_{k+1}, \ldots, \theta_N)
\end{align*}
where $n_k \coloneqq \lvert \theta_1 \rvert + \ldots + \lvert \theta_k \rvert
-k$.
A calculation similar to the one for the (reduced) cyclic complex shows that
$d$ and $b'$ are anti-commuting
differentials turning $\mathrm{B}(\Omega)$ into a $\Z_2$-graded complex with
differential $d + b'$ if we reduce the grading modulo two.
\subsubsection{Algebra valued bar cochains}
For a $\Z_2$-graded algebra $L$ the \textit{space of $L$-valued bar
cochains on $\Omega$} is defined as $\Hom(\mathrm{B}(\Omega), L)$, the space 
of linear maps from $\mathrm{B}(\Omega)$ to $L$.
By the universal property of the tensor product, every $L$-valued bar cochain
$l$ is given by a sequence $(l^{(N)})_{N \in \N_0}$ of multilinear maps
\begin{equation*}
	l^{(N)} \colon \underbrace{\Omega[1] \times \ldots \times \Omega[1]}_N
	\to L
\end{equation*}
where for $N = 0$ we may identify $l^{(0)}$ with an element of $L$.
We turn the space of $L$-valued bar cochains into a $\Z_2$-graded algebra with
product
\begin{equation}\label{Product bar cochains}
	l_1 l_2 (\theta_1, \ldots, \theta_N) \coloneqq \sum_{k=0}^N (-1)^{\lvert
		l_2 \rvert n_k} l_1(\theta_1, \ldots, \theta_k) l_2 (\theta_{k+1},
	\ldots,
	\theta_N)
\end{equation}
and $\Z_2$-grading given by the usual grading on $\Hom(V, W)$ for
$\Z_2$-graded
vector spaces $V$ and $W$.
The \textit{codifferential $\delta$} on $\Hom(\mathrm{B}(\Omega), L)$ is
defined via
\begin{equation*}
	\delta \coloneqq - (d+b')^{\vee} \colon \Hom(\mathrm{B}(\Omega), L) \to
	\Hom(\mathrm{B}(\Omega), L).
\end{equation*}
A standard calculation shows the following
\begin{lem}
	$\delta$ fulfills the graded Leibniz rule, i. e.
	\begin{equation}\label{Z2 graded Leibniz delta}
		\delta(l_1 l_2) = (\delta l_1)l_2 + (-1)^{\lvert l_1 \rvert} l_1
		(\delta
		l_2).
	\end{equation}
	In particular, $\Hom(\mathrm{B}(\Omega), L)$ is a differential
	$\Z_2$-graded algebra.
\end{lem}
\subsection{Analytic chains and entire cochains}\label{Section: entire}
Similarly to the previous section, we define the \textit{space of cyclic
complex valued cochains} to be the space $\Hom(\mathrm{C}(\Omega), \C)$ of
linear maps from $\mathrm{C}(\Omega)$ to $\C$.
To each cyclic cochain corresponds a
sequence $(l^{(N)})_{N \in \N_0}$ of multilinear maps
\begin{equation*}
	l^{(N)} \colon \Omega \times \underbrace{\underline{\Omega}[1] \times
		\ldots \times \underline{\Omega}[1]}_N \to \C.
\end{equation*}
\begin{defn}[\protect{\cite[Definition 3.1]{GüneysuLudewig}}]\label{Analytic}
	We call a cochain $l \in \Hom(\mathrm{C}(\Omega), \C)$ \textit{analytic}
	if there
	exists a continuous seminorm $\nu$ on $\Omega$ and a constant $C > 0$ such
	that for all $N \in \N_0$ and $\theta_0, \ldots, \theta_N \in \Omega$ we
	have
	\begin{equation*}
		\lvert l^{(N)}(\theta_0, \ldots, \theta_N) \rvert \leq
		\frac{C}{\lfloor N/2 \rfloor!} \, \nu(\theta_0) \cdot \ldots \cdot
		\nu(\theta_N).
	\end{equation*}
	We denote the space of analytic cochains by $\mathrm{C}_{\alpha}(\Omega)$.
\end{defn}
Recall that for a continuous seminorm $\nu$ on $\Omega$ the
\textit{$N$-th induced projective tensor seminorm $\pi_{\nu, N}$} on $\Omega
\otimes \underline{\Omega}[1]^{\otimes N}$ is given by
\begin{equation*}
	\pi_{\nu, N}(c_N) = \inf\left\{\sum_{\gamma} \nu(\theta_0^{(\gamma)})
	\cdot
	\ldots \cdot \nu(\theta_N^{(\gamma)}) : c_N = \sum_{\gamma}
	\theta_0^{(\gamma)} \otimes \ldots \otimes \theta_N^{(\gamma)}\right\},
\end{equation*}
the infimum being taken over all realizations of $c_N \in \mathrm{C}(\Omega)$
as a finite sum as in the above display.
If $l$ is an analytic cochain, it follows immediately from the definition
above that all induced linear maps $l^{(N)}$ are continuous on the projective 
tensor product $\Omega \otimes
\underline{\Omega}[1]^{\otimes N}$.
Moreover, since multiplication and the differential $d$ are assumed to be
continuous on $\Omega$, it follows that the codifferentials
$\underline{d}^\vee$,
$\underline{b}^\vee$ and $\underline{B}^\vee$ preserve the space
$\mathrm{C}_{\alpha}(\Omega)$.
\begin{defn}[\protect{\cite[Definition 3.2]{GüneysuLudewig}}]
	For a continuous seminorm $\nu$ on $\Omega$ the seminorm
	$\epsilon_\nu$ on $\mathrm{C}(\Omega)$ is given by
	\begin{equation}\label{Definition entire seminorm}
		\epsilon_{\nu}(c) \coloneqq \sum_{N = 0}^\infty \frac{\pi_{\nu,
				N}(c_N)}{\lfloor N/2 \rfloor!}, \qquad c = \sum_{N =
			0}^\infty c_N \in
		\mathrm{C}(\Omega), \quad \text{where $c_N \in \Omega \otimes
			\underline{\Omega}[1]^{\otimes N}$}.
	\end{equation}
	Note that this is well-defined since $c$ is a finite sum.
	Then, the \textit{space $\mathrm{C}^{\epsilon}(\Omega)$ of entire chains}
	is defined as the
	completion of $\mathrm{C}(\Omega)$ with respect to the seminorms
	$\epsilon_{\nu}$.
\end{defn}
If $l$ is an analytic cochain, it follows from the definitions that it is
continuous with respect to the seminorms $\epsilon_{\nu}$ and hence extends
to a continuous functional on $\mathrm{C}^\epsilon(\Omega)$.
Also, since $\underline{d}$, $\underline{b}$ and $\underline{B}$ are
continuous
on $\mathrm{C}(\Omega)$ with respect to the seminorms $\epsilon_\nu$, they
extend to differentials on $\mathrm{C}^\epsilon(\Omega)$ which are odd with
respect to the induced $\Z_2$-grading.
Hence, we have defined the \textit{entire (reduced) cyclic complex of
	$\Omega$}.
\subsection{Chen normalization}\label{Section: chen normalization}
Similarly to \cite{GetzlerJonesPetrack}, one often considers a quotient of the
reduced cyclic complex $\mathrm{C}(\Omega_{\mathbb{T}})$ by some subcomplex
$\mathrm{D}^\T(\Omega)$ since some of the chains which are of interest are
not closed in the larger complex.

To this end, let
\begin{equation*}
	S_i^{(f)} \colon \mathrm{C}(\Omega_{\mathbb{T}}) \to
	\mathrm{C}(\Omega_{\mathbb{T}})
\end{equation*}
be the continuous linear map given by
\begin{equation}\label{Si f}
	S_i^{(f)}(\theta_0, \dots, \theta_N) =
	\begin{cases}
		(-1)^{m_i}(\theta_0, \ldots, \theta_{i}, f, \theta_{i+1}, \ldots,
		\theta_N), & 0 \leq i \leq N - 1 \\
		(-1)^{m_N}(\theta_0, \ldots, \theta_N, f), & i = N \\
		0, & \text{otherwise},
	\end{cases}
\end{equation}
for all $f \in \Omega^0$ and $i\in \N_0$, and let
\begin{equation*}
	T_i^{(f)} \colon \mathrm{C}(\Omega_{\mathbb{T}}) \to
	\mathrm{C}(\Omega_{\mathbb{T}})
\end{equation*}
be the continuous linear map defined via
\begin{align*}
	T_i^{(f)}(\theta_0, \ldots, \theta_N) =
	\begin{cases}
		\begin{aligned}[b]
			& (\theta_0, \ldots, \theta_i, f \theta_{i+1}, \ldots, \theta_N)
			\\
			& \quad -(\theta_0, \ldots, \theta_i f, \theta_{i+1}, \ldots,
			\theta_N)
			\\
			& \quad -(\theta_0, \ldots, \theta_i, df, \theta_{i+1}, \ldots,
			\theta_N)
		\end{aligned}, & \quad 0 \leq i \leq N - 1 \\
		\begin{aligned}[b]
			& (f\theta_0, \ldots, \theta_N) \\
			& \quad -(\theta_0, \ldots, \theta_N f) \\
			& \quad -(\theta_0, \ldots, \theta_N, df)
		\end{aligned}, & \quad i = N
	\end{cases}
\end{align*}
for all $f \in \Omega^0$ and $i\in \N_0$.
Note that $\Omega^0 \subseteq \Omega_\T^0$ is a proper subset in general.
Moreover, define
\begin{align}
	&S \colon \mathrm{C}(\Omega_{\mathbb{T}}) \to
	\mathrm{C}(\Omega_{\mathbb{T}}), \quad (\theta_0, \dots, \theta_N) \mapsto
	\sum_{k=0}^N (\theta_0, \dots, \theta_{k}, \sigma, \theta_{k+1}, \dots,
	\theta_N),\label{Chen normalization S}\\
	&R \colon \mathrm{C}(\Omega_{\mathbb{T}}) \to
	\mathrm{C}(\Omega_{\mathbb{T}}), \quad (\theta_0, \dots, \theta_N) \mapsto
	(\sigma \theta_0, \dots, \theta_N).\label{Chen normalization R}
\end{align}
Finally, we let $\mathrm{D}^{\mathbb{T}}(\Omega) \subseteq
\mathrm{C}(\Omega_{\mathbb{T}})$ be the smallest subcomplex containing the
images of the operators
\begin{equation}\label{Chen operators}
	S + \mathbf{1}, \quad R, \quad S_{i}^{(f)}, \quad T_{i}^{(f)}
\end{equation}
for all $i \in \N_0$ and $f \in \Omega^0$.
\begin{defn}[\protect{\cite[Definition 3.4]{GüneysuLudewig}}]
	The quotient complex
	\begin{equation*}
		\mathrm{N}^{\T, \epsilon}(\Omega) \coloneqq
		\mathrm{C}^{\epsilon}(\Omega_\T) /
		\overline{\mathrm{D}^{\mathbb{T}}(\Omega)},
	\end{equation*}
	where the closure is taken inside $\mathrm{C}^\epsilon(\Omega_\T)$, is 
	called the \textit{Chen normalized entire complex}.
	The space $\mathrm{N}_{\T, \alpha}(\Omega)$ of analytic cochains 
	vanishing on $\mathrm{D}^{\mathbb{T}}(\Omega)$ will be referred to as the 
	\textit{space of extended Chen normalized analytic cochains}.
	This induces $\Z_2$-graded complexes
	\begin{equation*}
		\begin{tikzcd}
			\mathrm{N}_{+}^{\T, \epsilon}(\Omega) \ar[rr,  bend left=20,
			"\underline{d}+\underline{b}-\underline{B}"] & &  \ar[ll,  bend
			left=20, "\underline{d}+\underline{b}-\underline{B}"]
			\mathrm{N}_{-}^{\T, \epsilon}(\Omega),
		\end{tikzcd}
		~~~~\text{and}~~~~
		\begin{tikzcd}
			\mathrm{N}_{\T, \alpha}^+(\Omega) \ar[rr,  bend left=20,
			"(\underline{d}+\underline{b}-\underline{B})^{\vee}"] & &  \ar[ll,
			bend
			left=20, "(\underline{d}+\underline{b}-\underline{B})^{\vee}"]
			\mathrm{N}_{\T, \alpha}^-(\Omega).
		\end{tikzcd}
	\end{equation*}
\end{defn}
\section{\texorpdfstring{$\vartheta$}{ϑ}-summable
	\texorpdfstring{$\mathcal{C}_q$}{C(q)}-Fredholm modules}\label{Section:
	Cq FM}
For $q \in \N_0$ we denote by $\mathcal{C}_q$ the complex Clifford algebra
with generators $e_1, \ldots, e_q$ subject to the relations
\begin{equation*}
	e_i e_j + e_j e_i = -2 \delta_{ij}.
\end{equation*}
We turn $\mathcal{C}_q$ into a $*$-algebra by setting $e_i^* \coloneqq -e_i$
for
each $i$.
Moreover, $\mathcal{C}_q$ becomes a graded algebra by declaring the generators
$e_i$ to be odd.
Note that for $q = 0$ we have $\mathcal{C}_q = \C$.

Given a $\Z_2$-graded Hilbert space $\HH$ we denote by $\Lop(\HH)$ the
$\Z_2$-graded
$C^*$-algebra of bounded operators on $\HH$ where the involution is given by
the Hilbert space adjoint.
\begin{defn}
	A \textit{Hilbert module over $\mathcal{C}_q$} is a $\Z_2$-graded Hilbert
	space $\HH$
	together with a $\Z_2$-graded unital $*$-representation $\mathcal{C}_q \to
	\Lop(\HH)$.
\end{defn}
\begin{rem}
	\begin{enumerate}
		\item Note that for $q = 0$ this is nothing but a $\Z_2$-graded
		Hilbert space.
		\item If $\HH$ is a Hilbert module over $\mathcal{C}_q$, this
		canonically induces a Hilbert module structure on $\HH^m$ for $m \in
		\N$ by letting $\mathcal{C}_q$ act diagonally.
	\end{enumerate}
\end{rem}
\begin{nota}
	If $\HH$ is a Hilbert module over $\mathcal{C}_q$, we denote by
	$\Lop_{\mathcal{C}_q}(\HH)$ the subalgebra of $\Lop(\HH)$ consisting of
	operators on $\HH$ supercommuting with the $\mathcal{C}_q$-action.
	Furthermore, we define $\Lopclifqtrace(\HH)$ to be the subalgebra of
	trace-class operators supercommuting with the $\mathcal{C}_q$-action.
\end{nota}
\begin{defn}\label{Def FM}
	A \textit{$\vartheta$-summable $\mathcal{C}_q$-Fredholm module over a
	locally convex dg algebra $\Omega$} is a triple $\mathscr{M} = (\HH,
	\mathbf{c}, Q)$, where
	\begin{enumerate}[label=(\roman*)]
		\item $\HH$ is a Hilbert module over $\mathcal{C}_q$;
		\item $\mathbf{c} \colon \Omega \to \Lopclifq(\HH)$ is a
		continuous linear map which preserves the $\Z_2$-grading and 
		satisfies $\mathbf{c}(\mathbf{1}) = \mathbf{1}$,
		\item $Q$ is an odd, self-adjoint (possibly) unbounded operator on
		$\HH$ supercommuting with $\mathcal{C}_q$,
	\end{enumerate}
	such that the relations
	\begin{equation} \label{Multiplicativity}
		[Q, \mathbf{c}(f)] = \mathbf{c}(df), \quad \text{and} \quad
		\mathbf{c}(f\theta) = \mathbf{c}(f)\mathbf{c}(\theta), \quad
		\mathbf{c}(\theta f) = \mathbf{c}(\theta)\mathbf{c}(f)
	\end{equation}
	hold for all $f \in \Omega^0$ and $\theta \in \Omega$.
	Furthermore, the above data has to fulfill the following analytic 
	requirements:
	\begin{enumerate}
		\item[(A1)] For all $\theta\in \Omega$ the operators $C_\pm(\theta)
		\coloneqq \Delta^{\pm1/2}\mathbf{c}(\theta)\Delta^{\mp1/2}$ are
		densely defined and bounded, where $\Delta \coloneqq Q^2 + 
		\mathbf{1}$; moreover, the
		map $\theta \mapsto C_\pm(\theta)$ is continuous from $\Omega$
		to $\Lopclifq(\HH)$;
		\item[(A2)] For each $T>0$ one has $e^{-T Q^2}\in
		\Lopclifqtrace(\HH)$.
	\end{enumerate}
	If $\mathscr{M}$ satisfies all the above assumptions except
	\eqref{Multiplicativity}, we call $\mathscr{M}$ a
	\textit{$\vartheta$-summable weak
		$\mathcal{C}_q$-Fredholm module}.
\end{defn}
\begin{eg}
	For $q = 0$ a $\vartheta$-summable (weak) $\mathcal{C}_0$-Fredholm module
	is nothing else but an even $\vartheta$-summable (weak)
	Fredholm module in the sense of \cite[Definition 2.1]{GüneysuLudewig}.
\end{eg}
\begin{defn}\label{Acyclic extension FM}
	Let $\mathscr{M} = (\HH, \mathbf{c}, Q)$ be a weak $\vartheta$-summable
	$\mathcal{C}_q$-Fredholm module over a locally convex dg algebra $\Omega$.
	We extend $\mathscr{M}$ to a weak $\vartheta$-summable
	$\mathcal{C}_q$-Fredholm module $\mathscr{M}_\T = (\HH, \mathbf{c}_\T,
	Q)$ over $\Omega_\T$, the acyclic extension of $\Omega$ (\cref{Acyclic
		extension}), via
	\begin{equation*}
		\mathbf{c}_\T(\theta' + \sigma \theta'') \coloneqq
		\mathbf{c}(\theta').
	\end{equation*}
	We call this the \textit{acyclic extension of $\mathscr{M}$.}
\end{defn}
\begin{defn}
	An \textit{isomorphism} between two $\vartheta$-summable weak
	$\mathcal{C}_q$-Fredholm modules $\mathscr{M}_1 = (\HH_1, \mathbf{c}_1,
	Q_1)$ and $\mathscr{M}_2 = (\HH_2, \mathbf{c}_2, Q_2)$ over a locally
	convex dg algebra $\Omega$ consists of a unitary
	isomorphism $U \colon \HH_1 \to \HH_2$ of $\Z_2$-graded Hilbert spaces
	intertwining all structure maps, i. e.
	\begin{equation*}
		U(q v) = q (Uv) \quad \forall q \in \mathcal{C}_q, v \in \HH_1,
		\qquad
		\text{and} \qquad U(\mathbf{c}_1(\theta) v) =
		\mathbf{c}_2(\theta)(Uv)
		\quad \forall \theta \in \Omega, v \in \HH_1,
	\end{equation*}
	and $U (Q_1 v) = Q_2 (Uv)$ for all $v \in \mathrm{dom}(Q_1)$ with the
	implicit assumption that $U$ maps $\mathrm{dom}(Q_1)$ onto
	$\mathrm{dom}(Q_2)$.
\end{defn}
Finally, as in \cite[Definition 6.1]{GüneysuLudewig}, we introduce the notion
of a homotopy of $\vartheta$-summable $\mathcal{C}_q$-Fredholm modules:
\begin{defn}\label{Definition homotopy}
	A \textit{homotopy of $\vartheta$-summable (weak) $\mathcal{C}_q$-Fredholm
	modules over $\Omega$} is a family $\mathscr{M}^s = (\HH,
	\mathbf{c}^s,
	Q_s)$, $s \in [0,1]$, of $\vartheta$-summable (weak)
	$\mathcal{C}_q$-Fredholm modules over $\Omega$ (with fixed
	$\mathcal{C}_q$ action) such that:
	\begin{enumerate}
		\item[(H1)] In (A1) from \cref{Def FM} one can choose all seminorms 
		to be independent of $s$ and for all $T>0$, one has
		\begin{equation*}
			\sup_{s\in [0,1]}\Tr(e^{-TQ_s^2})<\infty.
		\end{equation*}
		\item[(H2)] The operators $Q_s$ are defined on the same domain for 
		all $s \in [0,1]$ and for all $h \in \mathrm{dom}(Q_s)$,
		the assignment $s \mapsto Q_s h$ is continuously differentiable from 
		$[0,1]$ to $\HH$.
		Thus, the time derivative $\dot{Q}_s$ is densely defined on $\HH$.
		The operators $\dot{Q}_s\Delta_s^{-1/2}$ and
		$\Delta_s^{-1/2}\dot{Q}_s$, where $\Delta_s 
		= Q_s^2+\mathbf{1}$, are assumed to be bounded in a uniform way:
		\begin{equation*}
			\sup_{s\in [0,1]}\bigl\|\Delta_s^{-1/2}\dot{Q}_s\bigr\| +
			\sup_{s\in
				[0,1]}\bigl\|\dot{Q}_s\Delta_s^{-1/2}\bigr\|<\infty.
		\end{equation*}
		\item[(H3)] The map $s\mapsto
		\mathbf{c}^{s}(\theta)$ defines a continuously differentiable curve 
		for all $\theta\in \Omega$ if we equip $\Lop(\HH)$ with the strong 
		operator topology.
		Moreover, for each $\theta\in \Omega$ and $s \in [0,1]$, the
		operators $\dot{C}^s_\pm(\theta) \coloneqq
		\Delta_s^{\pm1/2}\dot{\mathbf{c}}^s(\theta)\Delta_s^{\mp1/2}$ are
		densely defined and bounded and the assignment $\theta \mapsto
		\dot{C}^s_\pm(\theta)$ is continuous from $\Omega$ to
		$\Lopclifq(\HH)$ in such a way that the seminorms can be chosen
		independently of $s$.
	\end{enumerate}
\end{defn}
\section{Odd \texorpdfstring{$\vartheta$}{ϑ}-summable Fredholm
	modules}\label{Section: odd FM}
In this section, we will introduce the notion of an \textit{odd}
$\vartheta$-summable Fredholm module over a locally convex dg algebra and see
that this induces a $\vartheta$-summable $\mathcal{C}_1$-Fredholm module in a
canonical way.
\begin{defn}\label{Definition odd FM}
	An \textit{odd $\vartheta$-summable Fredholm module} $\mathscr{M}$ over a
	locally convex dg algebra $\Omega$ is
	a triple $(\HH, \mathbf{c}, Q)$, where
	\begin{enumerate}
		\item $\HH$ is a Hilbert space;
		\item $\mathbf{c} \colon \Omega \to \Lop(\HH)$ is a continuous, linear
		operator such that $\mathbf{c}(\mathbf{1}) = \mathbf{1}$;
		\item $Q$ is an unbounded, self-adjoint operator on $\HH$,
	\end{enumerate}
	such that the following holds for all $f \in \Omega^0$ and $\theta \in
	\Omega$:
	\begin{equation}\label{Equations Fredholm}
		[Q, \mathbf{c}(f)] = \mathbf{c}(df), \quad \text{and} \quad
		\mathbf{c}(f \theta) =
		\mathbf{c}(f) \mathbf{c}(\theta), \quad \mathbf{c}(\theta f) =
		\mathbf{c}(\theta) \mathbf{c}(f).
	\end{equation}
	Furthermore, the above data are subject to analytic requirements:
	\begin{enumerate}[label=(A{{\arabic*}})]
		\item For all $\theta \in \Omega$, the operators $C_{\pm}(\theta)
		\coloneqq \Delta^{\pm 1/2} \mathbf{c}(\theta) \Delta^{\mp 1/2}$ are
		densely defined and bounded where $\Delta \coloneqq Q^2 + \mathbf{1}$
		and the map $\theta \mapsto C_{\pm}(\theta)$ is continuous from
		$\Omega$ to $\Lop(\HH)$;
		\item For all $T > 0$, one has $e^{-TQ^2} \in \Lop^1(\HH)$.
	\end{enumerate}
	If $\mathscr{M}$ satisfies all properties apart from \cref{Equations
	Fredholm}, it is called an \textit{odd $\vartheta$-summable weak Fredholm
	module over $\Omega$}.
\end{defn}
Suppose we are given an odd $\vartheta$-summable Fredholm module
$\mathscr{M}$.
Then, we may define a $\vartheta$-summable $\mathcal{C}_1$-Fredholm module
$\tilde{\mathscr{M}} \coloneqq (\tilde{\HH}, \tilde{\mathbf{c}},
\tilde{Q})$ as follows:
\begin{enumerate}
	\item Let $\tilde{\HH} \coloneqq \HH \oplus \HH$ with
	$\Z_2$-grading determined by $\tilde{\HH}^+ \coloneqq \HH \oplus
	\{0\}$,
	$\tilde{\HH}^- \coloneqq \{0\} \oplus \HH$;
	\item Let $\mathcal{C}_1$ act on $\tilde{\HH}$ by letting the
	generator
	$e_1$ act on $\tilde{\HH}$ via
	\begin{equation*}
		\gamma \coloneqq
		\begin{pmatrix}
			0 & \mathbf{1} \\
			-\mathbf{1} & 0
		\end{pmatrix};
	\end{equation*}
	\item
	\begin{equation*}
		\tilde{\mathbf{c}}(\theta) \coloneqq
		\begin{cases}
			\begin{pmatrix}
				\mathbf{c}(\theta) & 0 \\
				0 & \mathbf{c}(\theta)
			\end{pmatrix}, & \text{$\theta$ even} \\[15pt]
			\begin{pmatrix}
				0 & \mathbf{c}(\theta) \\
				\mathbf{c}(\theta) &0
			\end{pmatrix}, & \text{$\theta$ odd}
		\end{cases} \qquad \text{for $\theta \in \Omega$};
	\end{equation*}
	\item
	\begin{equation*}
		\tilde{Q} \coloneqq
		\begin{pmatrix}
			0 & Q \\
			Q & 0
		\end{pmatrix}.
	\end{equation*}
\end{enumerate}
\begin{rem}
	If $\mathscr{M}$ is an odd $\vartheta$-summable \textit{weak} Fredholm
	module, then $\tilde{\mathscr{M}}$ will be an even $\vartheta$-summable
	\textit{weak} Fredholm module.
\end{rem}
\begin{rem}\label{Doubling acyclic}
	Consider the acyclic extension $\Omega_{\mathbb{T}}$ of
	$\Omega$.
	As in \cref{Acyclic extension FM}, we may define the extended odd
	$\vartheta$-summable weak Fredholm module $\mathscr{M}_{\mathbb{T}} =
	(\HH,\mathbf{c}_{\mathbb{T}}, Q)$ over $\Omega_{\mathbb{T}}$ by setting
	\begin{equation*}
		\mathbf{c}_{\mathbb{T}}(\theta' + \sigma \theta'') \coloneqq
		\mathbf{c}(\theta').
	\end{equation*}
	In view of the preceding example, it is obvious that
	$\widetilde{\mathscr{M}_{\mathbb{T}}} =
	\tilde{\mathscr{M}}_{\mathbb{T}}$.
\end{rem}
\section{The Chern character}\label{Section: Chern character}
Similarly to \cite{GüneysuLudewig}, for a $\vartheta$-summable weak
$\mathcal{C}_q$-Fredholm module $\mathscr{M}$, we will construct an analytic
cochain $\Chern$ which we will call the \textit{Chern character of
	$\mathscr{M}$}.
\subsection{The quantization map}
Let $\mathscr{M} \coloneqq (\mathcal{H}, \mathbf{c}, Q)$ be a
$\vartheta$-summable weak $\mathcal{C}_q$-Fredholm module.
Similarly to \cite[Equation (4.1)]{GüneysuLudewig}, we define an
$\Lopclifq(\HH)$-valued bar cochain $\omega_{\mathscr{M}}$ via
\begin{equation*}
	\omega_{\mathscr{M}}^{(0)} \coloneqq Q, \quad
	\omega_{\mathscr{M}}^{(1)}(\theta) \coloneqq \mathbf{c}(\theta), \quad
	\omega_{\mathscr{M}}^{(N)}(\theta_1, \ldots, \theta_N) \coloneqq 0, N \geq
	2.
\end{equation*}
The cochain $\omega_{\mathscr{M}}$ is odd in the $\Z_2$-grading of
$\Hom(\mathrm{B}(\Omega), \Lop_{\mathcal{C}_q}(\HH))$.
Note that since $Q$ may be unbounded in general (hence not actually an
element of
$\Lop_{\mathcal{C}_q}(\HH)$), the above display is only formal.
We may interpret $\omega_{\mathscr{M}}$ as a (super-)connection form of the
(super-)connection $\nabla \coloneqq \delta \omega_{\mathscr{M}} +
\omega_{\mathscr{M}}$ and define its curvature via
\begin{equation*}
	F_{\mathscr{M}} \coloneqq \delta \omega_{\mathscr{M}} +
	\omega_{\mathscr{M}}^2.
\end{equation*}

Working out the components, $F_{\mathscr{M}}$ is explicitly given by
\begin{align}\label{Components curvature}
	F_{\mathscr{M}}^{(0)} & = Q^2, \notag\\
	F_{\mathscr{M}}^{(1)}(\theta_1) & = [Q, \mathbf{c}(\theta_1)] - \mathbf{c}(d
	\theta_1), \\
	F_{\mathscr{M}}^{(2)}(\theta_1, \theta_2) & = (-1)^{\lvert
		\theta_1 \rvert}(\mathbf{c}(\theta_1 \theta_2) - \mathbf{c}(\theta_1)
	\mathbf{c}(\theta_2)),\notag
\end{align}
with all higher arities being zero.
\begin{nota}
	Let $H$ be a non-negative operator and $A_1, \ldots, A_N$ suitable
	operators on a Hilbert space $\HH$.
	Define
	\begin{equation*}
		\{A_1, \ldots, A_N\}_{H} \coloneqq \int_{\Delta_N} e^{-\tau_1 H} A_1
		e^{-(\tau_2 - \tau_1)H} A_2 \cdot \ldots \cdot e^{-(\tau_{N} -
			\tau_{N-1})H}A_n e^{-(1 - \tau_N)H} \, \mathrm{d} \tau
	\end{equation*}
	where $\Delta_N\coloneqq \{\tau \in \R^N \mid 0 \leq \tau_1 \leq \ldots
	\leq \tau_N \leq 1\}$ is the standard simplex.
\end{nota}
\cite[Lemma 4.2]{GüneysuLudewig} shows that the components of the curvature
$F_{\mathscr{M}}$ satisfy the necessary analytic assumptions such that the
following definition of the \textit{quantization map} makes sense:
\begin{defn}\label{Defn Quantization}
	For $T>0$ the even $\Lop_{\mathcal{C}_q}(\HH)$-valued bar cochain 
	$\Phi^{\mathscr{M}}_T$ is defined as
	\begin{equation*}
		\Phi^{\mathscr{M}}_T \coloneqq \sum_{N=0}^\infty (-T)^N
		\bigl\{\underbrace{F^{\geq 1}_{\mathscr{M}}, \dots, F^{\geq
		1}_{\mathscr{M}}}_N\bigr\}_{TQ^2}.
	\end{equation*}
\end{defn}
We recall:
\begin{nota}
	Let $\mathcal{P}_{M, N}$ be the set of ordered partitions of length $M$ 
	of $\{1, \ldots, N\}$; that is, $M$-tuples of nonempty subsets
	$I = (I_1, \ldots, I_M)$ for which $I_1 \cup \ldots \cup I_M = \{1, \dots,
	N\}$ and for which every member of $I_a$ is smaller than any member of
	$I_b$ if $a < b$.
\end{nota}
For $N \geq 1$, we have the formula (see \cite[Eq.
(4.8)]{GüneysuLudewig})
\begin{equation*}
	\Phi^{\mathscr{M}}_T(\theta_1, \dots, \theta_N) = \sum_{M=1}^N (-T)^M
	\sum_{I \in
		\mathscr{P}_{M, N}} \bigl\{ F_{\mathscr{M}}(\theta_{I_1}), \dots,
	F_{\mathscr{M}}(\theta_{I_M})\bigr\}_{TQ^2},
\end{equation*}
where $\theta_{I_a}:= (\theta_{i+1} , \dots , \theta_{i+m})$
if $I_a = \{j \mid i < j \leq i+m\}$ for some $i, m$.
For $N = 0$ we have $(\Phi_T^{\mathscr{M}})^{(0)} = e^{-TQ^2}$.
\begin{rem}\label{Remark isomorphism quantm}
	\begin{enumerate}
		\item We have that $\Phi \in \Hom^+(\mathrm{B}(\Omega),
		\Lop_{\mathcal{C}_q}(\HH))$, i. e. it maps even elements of
		$\mathrm{B}(\Omega)$ to even operators on $\HH$ and odd elements to
		odd operators, since $F_{\mathscr{M}}$ is even (because $\omega$ is
		odd).
		\item\label{Item 3} If $U \colon \mathscr{M}_1 \to \mathscr{M}_2$ is
		an isomorphism
		of $\vartheta$-summable weak $\mathcal{C}_q$-Fredholm modules, it is
		easy to check that
		\begin{equation*}
			U \Phi_T^{\mathscr{M}_1}(\theta) = \Phi_T^{\mathscr{M}_2}(\theta)
			U
		\end{equation*}
		for all $\theta \in \mathrm{B}(\Omega)$.
	\end{enumerate}
\end{rem}
As in \cite[Theorem 4.6]{GüneysuLudewig} one proves the following fundamental
estimate:
\begin{thm}\label{Fundamental estimate}
	The operator $\Phi_{T}^{\mathscr{M}}(\theta_1, \ldots, \theta_N)$ is 
	well-defined and trace-class for all $T > 0$ and $\theta_1, \ldots, 
	\theta_N \in \Omega$.
	There exists a continuous seminorm $\nu$ on $\Omega$, which is independent
	of $T$, such that
	\begin{equation*}
		\lVert \Phi_{T}^{\mathscr{M}}(\theta_1, \ldots, \theta_N) \rVert_1
		\leq
		e^{T/2} \Tr(e^{-TQ^2/2}) \frac{T^{N/2}}{\lfloor N/2 \rfloor!}
		\nu(\theta_1) \cdot \ldots \cdot \nu(\theta_N).
	\end{equation*}
	Moreover, the same holds for $Q \Phi_{T}^{\mathscr{M}}$ and
	$\Phi_T^{\mathscr{M}} Q$ instead of $\Phi_T^{\mathscr{M}}$ if we replace
	$T^{N/2}$ with $T^{(N-1)/2}$ and $\lfloor N/2\rfloor!$ with $\lfloor
	(N-1)/2 \rfloor!$ on the right-hand side.
\end{thm}
\subsection{Definition of \texorpdfstring{$\Chern$}{ChernM}}
\begin{defn}
	Let $\HH$ be a Hilbert module over $\mathcal{C}_q$ and $A \in
	\Lopclifqtrace(\HH)$.
	The \textit{Clifford supertrace of $A$} is defined as
	\begin{equation*}
		\CStr(A) \coloneqq 2^{-q}\Str(\gamma A)
	\end{equation*}
	where $\gamma \coloneqq e_1 \cdot \ldots \cdot e_q$ is the Clifford volume
	element.
\end{defn}
The following lemma justifies the name \enquote{Clifford supertrace}:
\begin{lem}\label{Clifford trace is trace}
	Let $\HH$ be a Hilbert module over $\mathcal{C}_q$ and $A, B \in
	\Lopclifqtrace(\HH)$.
	Then one has
	\begin{equation*}
		\CStr([A, B]) = 0.
	\end{equation*}
\end{lem}
With the above definitions in place, we may define the Chern character:
\begin{defn}
	Let $\mathscr{M}$ be a $\vartheta$-summable weak $\mathcal{C}_q$-Fredholm
	module.
	The \textit{Chern character of $\mathscr{M}$} is the cochain $\Chern \in
	\Hom(\mathrm{C}(\Omega), \C)$ defined via
	\begin{equation*}\label{Definition Chern character of FM}
		\Chern(\theta_0, \ldots, \theta_N) \coloneqq
		\CStr(\mathbf{c}(\theta_0)
		\Phi_1^{\mathscr{M}}(\theta_1, \ldots, \theta_N))
	\end{equation*}
	for $\theta_0, \ldots, \theta_N \in \Omega$.
\end{defn}
\begin{rem}
	For $q = 0$ this coincides with \cite[Definition 5.1]{GüneysuLudewig}.
\end{rem}
\begin{lem}\label{Chern even/odd}
	$\Chern \in \Hom^+(\mathrm{C}(\Omega), \C)$ for $q$ even and $\Chern \in
	\Hom^-(\mathrm{C}(\Omega), \C)$ for $q$ odd.
\end{lem}
\begin{proof}
	For $q$ even we have to prove that $\Chern$ is even, i. e. it vanishes on
	odd elements of $\mathrm{C}(\Omega)$.
	Let $(\theta_0, \ldots, \theta_N) \in \mathrm{C}(\Omega)$ be an odd
	element.
	There are two cases to consider:
	\begin{enumerate}
		\item $\theta_0$ is even in $\Omega$.
		By definition of the gradings on $\mathrm{C}(\Omega)$ and
		$\mathrm{B}(\Omega)$ this implies that $(\theta_1, \ldots, \theta_N)$
		is odd in $\mathrm{B}(\Omega)$.
		Since $\mathbf{c}$ and $\Phi_1^{\mathscr{M}}$ are parity preserving,
		it
		follows that
		\begin{equation*}
			\mathbf{c}(\theta_0) \Phi_1^{\mathscr{M}}(\theta_1, \ldots,
			\theta_N)
		\end{equation*}
		is an odd operator on $\HH$.
		Thus, as $\gamma$ is even (because $q$ is), $\gamma
		\mathbf{c}(\theta_0)\Phi_1^{\mathscr{M}}(\theta_1, \ldots, \theta_N)$
		is an odd operator on which the usual supertrace on $\HH$ vanishes.
		\item $\theta_0$ is odd in $\Omega$.
		In this case, the same argument shows that $(\theta_1, \ldots,
		\theta_N)$ must be even in $\mathrm{B}(\Omega)$ and the above proof
		goes through in a similar fashion.
	\end{enumerate}
	The case in which $q$ is odd may be proven analogously.
\end{proof}
\begin{lem}\label{Isomorphism Chern character}
	If $\mathscr{M}_1$ and $\mathscr{M}_2$ are isomorphic
	$\vartheta$-summable weak $\mathcal{C}_q$-Fredholm modules, then
	$\mathrm{Ch}_{\mathscr{M}_1} = \mathrm{Ch}_{\mathscr{M}_2}$.
\end{lem}
\begin{proof}
	This follows from \crefdefpart{Remark isomorphism quantm}{Item
		3} and the fact that the isomorphism intertwines the
	$\mathcal{C}_q$-actions.
\end{proof}
We collect some of the most important properties of the Chern character:
\begin{thm}\label{Thm ChM}
	Let $\mathscr{M}$ be a $\vartheta$-summable weak $\mathcal{C}_q$-Fredholm
	module.
	Then the following assertions hold:
	\begin{enumerate}
		\item The cochain $\Chern$ is analytic.
		\item The Chern character is coclosed, that is $(\underline{d} +
		\underline{b} - \underline{B})^\vee \Chern = 0$.
		\item If $\mathscr{M}_\T$ denotes the acyclic extension of
		$\mathscr{M}$, then $\ChernT$ is Chen normalized, i. e. it vanishes
		on the subcomplex $\mathrm{D}^{\mathbb{T}}(\Omega)$ defined in
		\cref{Section: chen normalization}.
		\item For $\mathscr{M}_s$, $s \in [0,1]$, a homotopy of
		$\vartheta$-summable $\mathcal{C}_q$-Fredholm modules
		the Chern characters $\mathrm{Ch}_{\mathscr{M}^0_\T}$ and
		$\mathrm{Ch}_{\mathscr{M}^1_\T}$ are homologous in
		$\mathrm{N}_{\T, \alpha}(\Omega)$.
	\end{enumerate}
\end{thm}
\subsubsection{Proof of \texorpdfstring{\cref{Thm ChM}}{Theorem 5.10}, Part 1}
Analyticity of $\Chern$ follows from the fundamental estimate in
\cref{Fundamental estimate} together with the fact that
\begin{equation*}
	\lVert AB \rVert_1 \leq \lVert A \rVert \cdot \lVert B \rVert_1
\end{equation*}
for bounded operators $A$ and $B$ with $B$ trace-class.\qed
\subsubsection{Proof of \texorpdfstring{\cref{Thm ChM}}{Theorem 5.10}, Part 2}
We will adapt the proof of \cite[Theorem 5.3]{GüneysuLudewig} to
our more general setting and highlight only the main differences.

The inclusion $\Omega \subseteq \Omega_{\mathbb{T}}$ gives rise to a linear 
map
\begin{equation*}
	j \colon \mathrm{C}(\Omega) \to \mathrm{C}(\Omega_{\mathbb{T}})
\end{equation*}
which commutes with all differentials since $\Omega \subseteq 
\Omega_{\mathbb{T}}$ is a subalgebra and $d_{\mathbb{T}}$ restricts to $d$ on 
$\Omega$.
Observe that by definition of the acyclic extension of $\mathscr{M}$, we have
$\Chern = j^*\ChernT$ so that it suffices to prove coclosedness of $\ChernT$.

The maps $\Phi_1^{\mathscr{M}_{\mathbb{T}}}$ and
$F_{\mathscr{M}_{\mathbb{T}}}$ descend to the quotient
$\mathrm{B}(\underline{\Omega_{\mathbb{T}}})$ which is defined as in \cref{Bar
construction} with the graded vector space $\underline{\Omega_{\mathbb{T}}}$
instead of the dg algebra $\Omega_{\mathbb{T}}$; the quotient map of
$\Phi_1^{\mathscr{M}_{\mathbb{T}}}$ will be denoted by
$\underline{\Phi}_1^{\mathscr{M}_{\mathbb{T}}}$.
Using the explicit formulas from \cref{Components curvature} we
have
\begin{align*}
	F_{\mathscr{M}_{\mathbb{T}}}^{(0)} & = Q^2, \\
	F_{\mathscr{M}_{\mathbb{T}}}^{(1)}(\theta_1) & =
	[Q, \mathbf{c}(\theta_1')] - \mathbf{c}(d
	\theta_1') + \mathbf{c}(\theta''), \\
	F_{\mathscr{M}_{\mathbb{T}}}^{(2)}(\theta_1, \theta_2) & =
	(-1)^{\lvert
		\theta_1 \rvert}(\mathbf{c}(\theta_1' \theta_2') -
	\mathbf{c}(\theta_1')\mathbf{c}(\theta_2'))
\end{align*}
for $\theta_1, \theta_2 \in \Omega_{\mathbb{T}}$.
Let us introduce some notation:
define
\begin{equation*}
	\alpha \colon \mathrm{C}(\Omega_{\mathbb{T}}) \to
	\mathrm{B}(\underline{\Omega_{\mathbb{T}}}), \quad (\theta_0, \ldots,
	\theta_N) \mapsto \underline{\mathbf{N}}(\sigma \theta_0, \theta_1,
	\ldots,
	\theta_N)
\end{equation*}
to be the parity preserving map where
\begin{equation}\label{Definition alpha}
	\underline{\mathbf{N}} \colon \mathrm{B}(\underline{\Omega_{\mathbb{T}}})
	\to \mathrm{B}(\underline{\Omega_{\mathbb{T}}})
\end{equation}
is the quotient map of
\begin{equation*}
	\mathbf{N} \colon \mathrm{B}(\Omega_{\mathbb{T}}) \to
	\mathrm{B}(\Omega_{\mathbb{T}}), \quad (\theta_1, \ldots, \theta_N)
	\mapsto
	\sum_{k = 1}^N (-1)^{n_k(n_N - n_k)} (\theta_{k+1}, \ldots, \theta_N,
	\theta_1, \ldots, \theta_k)
\end{equation*}
and $n_k = \lvert \theta_1 \rvert + \ldots + \lvert \theta_k \rvert - k$.
\begin{lem}\label{Ch pullback}
	The Chern character of the acyclic extension satisfies
	\begin{equation*}
		\mathrm{Ch}_{\mathscr{M}_{\mathbb{T}}} = - \alpha^*
		\CStr(\underline{\Phi}_1^{\mathscr{M}_{\mathbb{T}}}).
	\end{equation*}
\end{lem}
\begin{proof}
	If $N = 0$, the statement is obvious.
	For $N \geq 1$, we calculate
	\begin{align*}
		& \alpha^*
		\CStr(\underline{\Phi}_1^{\mathscr{M}_{\mathbb{T}}})
		(\theta_0, \ldots, \theta_N) \\
		& \qquad = \sum_{k = 1}^{N+1} (-1)^{m_{k-1}(m_N - m_{k-1})}\CStr(
		\underline{\Phi}_1^{\mathscr{M}_{\mathbb{T}}}(\theta_k, \ldots,
		\theta_N, \sigma \theta_0, \theta_1, \ldots, \theta_{k-1}))
	\end{align*}
	where $m_k = \lvert \theta_0 \rvert + \ldots +
	\lvert \theta_k \rvert - k$.
	Indeed, the signs in the above equation follow from
	\begin{equation*}
		n_k(\sigma \theta_0, \ldots, \theta_N) = \lvert \sigma \theta_0 \rvert
		+ \lvert \theta_1 \rvert + \ldots + \lvert \theta_{k-1} \rvert - k =
		m_{k-1} \mod 2.
	\end{equation*}
	Using the explicit formulas for the curvature, one notices that
	\begin{equation*}
		F_{\mathscr{M}_{\mathbb{T}}}(\theta_N, \sigma\theta_0) =
		F_{\mathscr{M}_{\mathbb{T}}}(\sigma \theta_0, \theta_1) = 0
		\qquad \text{and} \qquad
		F_{\mathscr{M}_{\mathbb{T}}}(\sigma\theta_0) = \mathbf{c}(\theta_0').
	\end{equation*}
	Hence, expanding $\underline{\Phi}_1^{\mathscr{M}_{\mathbb{T}}}$ in
	terms of the brackets yields
	\begin{equation*}
		\sum_{k=1}^{N+1} (-1)^{m_{k-1}(m_N-m_{k-1})} \sum_{M=1}^N (-1)^{M+1}
		\sum_{l=1}^{M+1} \sum_{I \in \mathcal{P}_{M, N}^{k, l}}
		\begin{aligned}
			\CStr &\bigl\{F_{\mathscr{M}_{\mathbb{T}}}
			(\theta_{I_{l}}), \dots \\
			& \dots, F_{\mathscr{M}_{\mathbb{T}}}(\theta_{I_M}),
			\mathbf{c}(\theta_0^\prime),
			F_{\mathscr{M}_{\mathbb{T}}}(\theta_{I_1}), \dots \\
			&~~~~~~~~~~~~~~~~~~~~~\dots,
			F_{\mathscr{M}_{\mathbb{T}}}(\theta_{I_{l-1}})
			\bigr\}_{Q^2},
		\end{aligned}
	\end{equation*}
	where $\mathcal{P}^{k, l}_{M, N} \subseteq \mathcal{P}_{M, N}$ are those 
	partitions for which $I_1, \dots, I_{\ell-1} \subseteq \{1, \dots,
	k-1\}$ and $I_{l}, \dots, I_M \subseteq \{k, \dots, N\}$.
	Let $I \in \mathcal{P}^{k, l}_{M, N}$ such that the bracket in the above
	display is non-zero and let $1 \leq j \leq l-1$.
	On the one hand, if $\theta_{I_j} = (\theta_{i_j}, \theta_{i_{j}+1})$, we
	have
	\begin{align*}
		\lvert F_{\mathscr{M}_{\mathbb{T}}}(\theta_{I_j}) \rvert & =
		\lvert \mathbf{c}(\theta_{i_j}' \theta_{i_j + 1}') -
		\mathbf{c}(\theta_{i_j}')\mathbf{c}(\theta_{i_j + 1}') \rvert =
		\lvert \theta_{i_j}' \rvert + \lvert \theta_{i_j + 1}' \rvert \\
		& = \lvert \theta_{i_j} \rvert + \lvert \theta_{i_j + 1} \rvert - 2
		\quad \mod 2,
	\end{align*}
	while on the other hand, if $\theta_{I_j} = \theta_{i_j}$, we have
	\begin{equation*}
		\lvert F_{\mathscr{M}_{\mathbb{T}}}(\theta_{I_J}) \rvert =
		\lvert [Q, \mathbf{c}(\theta_{i_j}')] -
		\mathbf{c}(d
		\theta_{i_j}') + \mathbf{c}(\theta_{i_j}'') \rvert = \lvert
		\theta_{i_j}' \rvert + 1 = \lvert \theta_{i_j} \rvert - 1 \mod 2.
	\end{equation*}
	Thus, we have
	\begin{equation*}
		m_{k-1} = \lvert \theta_0' \rvert + \lvert
		F_{\mathscr{M}_{\mathbb{T}}}(\theta_{I_1}) \rvert + \ldots +
		\lvert F_{\mathscr{M}_{\mathbb{T}}}(\theta_{I_{l-1}}) \rvert
		\eqqcolon \tilde{m}_{l-1}
	\end{equation*}
	modulo two for $1 \leq k \leq N + 1$.
	Note that for all $l$ and $M \leq N$, the sets $\mathcal{P}_{M,
	N}^{k, l}$ and $\mathcal{P}_{M,
	N}^{k', l}$ are disjoint for $k \neq k'$ and their union over $k$ is  
	$\mathcal{P}_{M, N}$.
	Thus, the above sums over $k$ and $\mathcal{P}_{M, N}^{k, l}$ may be 
	written as one sum over $\mathcal{P}_{M, N}$, yielding
	\begin{equation}\label{Chern pullback long sum}
		- \sum_{M=1}^N (-1)^M \sum_{I \in \mathcal{P}_{M, N}} \sum_{l=1}^{M+1}
		(-1)^{\tilde{m}_{l-1}(\tilde{m}_M - \tilde{m}_{l-1})}
		\begin{aligned}
			\CStr
			&\bigl\{F_{\mathscr{M}_{\mathbb{T}}}(\theta_{I_{l}}), \dots
			\\
			& \dots, F_{\mathscr{M}_{\mathbb{T}}}(\theta_{I_M}),
			\mathbf{c}(\theta_0^\prime),
			F_{\mathscr{M}_{\mathbb{T}}}(\theta_{I_1}), \dots \\
			&~~~~~~~~~~~~~~~~~~~~~\dots,
			F_{\mathscr{M}_{\mathbb{T}}}(\theta_{I_{l-1}})
			\bigr\}_{Q^2}.
		\end{aligned}
	\end{equation}
	To simplify the above sum, we need the following
	\begin{lem}\label{Clifford trace cyclic sum property}
		Let $\HH$ be a Hilbert module over $\mathcal{C}_q$, $H$ a
		non-negative self-adjoint even
		operator on $\HH$ and $A_0, \ldots, A_N$ as in \cite[Lemma
		4.2]{GüneysuLudewig} with
		the extra condition that $H$ and the $A_j$ supercommute with the
		Clifford action.
		Then, the Clifford supertrace has the following cyclic property:
		\begin{equation*}
			\sum_{j = 0}^N (-1)^{k_j}\CStr(\{A_{j+1}, \ldots, A_N, A_0,
			\ldots,
			A_j\}_H) =
			\CStr(A_0 \{A_1, \ldots, A_N\}_H)
		\end{equation*}
		where $k_j \coloneqq (\lvert A_0 \rvert + \ldots + \lvert A_j
		\rvert)(\lvert A_{j+1} \rvert + \ldots + \lvert A_N \rvert)$.
	\end{lem}
	\begin{proof}
		As a consequence of \cref{Clifford trace is trace} we have
		\begin{align*}
			& \CStr\big([A_0 e^{-\tau_1 H}A_1 e^{-(\tau_2 - \tau_1)H} \cdot
			\ldots
			\cdot
			e^{-(\tau_j
				- \tau_{j-1})H} A_j e^{-(\tau_{j+1} - \tau_j)H}, \\
			& \qquad A_{j+1} e^{-(\tau_{j+2} - \tau_{j+1}) H}A_{j+2}
			e^{-(\tau_{j+3} - \tau_{j+2})H} \cdot \ldots \cdot e^{-(\tau_N
				- \tau_{N-1})H} A_N e^{-(1-\tau_N)H}]\big) = 0.
		\end{align*}
		Thus, after a substitution in $\Delta_N$, we have
		\begin{equation*}
			\CStr(A_0 \{A_1, \ldots, A_N\}_H) = (-1)^{k_j}\CStr(A_{j+1}
			\{A_{j+2},
			\ldots,
			A_N, A_0, \ldots, A_j\}_H)
		\end{equation*}
		for any $0 \leq j \leq N-1$ and in particular,
		\begin{equation}\label{Trace cyclic with 1}
			\CStr(A_0 \{A_1, \ldots, A_j, \mathbf{1}, A_{j+1}, \ldots,
			A_N\}_H) =
			(-1)^{k_j}\CStr(\{A_{j+1}, \ldots, A_N,
			A_0, \ldots, A_j\}_H).
		\end{equation}
		\begin{claim}
			For $0 \leq j \leq N$ we have
			\begin{equation*}
				\{A_1, \ldots, A_j, \mathbf{1}, A_{j+1}, \ldots, A_N\}_H =
				\int_{\Delta_N} (\tau_{j+1} - \tau_j) e^{-\tau_1 H}
				\prod_{j=1}^{N}
				A_j e^{-(\tau_{j+1} - \tau_j)H} \, \mathrm{d} \tau
			\end{equation*}
			where $\tau_{N+1} = 1$ and $\tau_0 = 0$.
		\end{claim}
		\begin{proof}[Proof of Claim]
			This follows from exploiting the semigroup property together with
			integration by parts:
			\footnotesize
			\begin{alignat*}{1}
				& \{A_1, \ldots, A_j, \mathbf{1}, A_{j+1}, \ldots, A_N\}_H \\
				& \qquad =
				\int_{\Delta_{N+1}} e^{-\tau_1 H}A_1
				\cdot \ldots \cdot A_j e^{-(\tau_{j+1} - \tau_j)H}
				\mathbf{1}
				e^{-(\tau_{j+2} - \tau_{j+1})H} A_{j+1} \cdot \ldots \cdot
				A_{N+1} e^{-(1 -
					\tau_{N+1})H} \, \dint \tau \\
				& \qquad =
				\int_{\Delta_{N+1}} e^{-\tau_1 H}A_1
				\cdot \ldots \cdot A_j
				e^{-(\tau_{j+2} - \tau_{j})H} A_{j+1} \cdot \ldots \cdot
				A_{N+1}
				e^{-(1 - \tau_{N+1})H} \, \dint \tau \\
				& \qquad = \int_0^1 \int_{\tau_1}^{1} \ldots
				\int_{\tau_{j-1}}^{1}
				\int_{\tau_{j+1}}^{1} \ldots \int_{\tau_{N}}^{1} \ldots \\
				& \qquad ~~~~~~~~~~~~~ \ldots
				\int_{\tau_{j-1}}^{\tau_{j+1}} \ldots \cdot A_j
				e^{-(\tau_{j+2} - \tau_{j})H} A_{j+1} \cdot \ldots \, \dint
				\tau_j \, \dint \tau_{N+1} \, \ldots
				\dint \tau_{j+2} \, \dint \tau_{j+1} \, \ldots \, \dint
				\tau_1.
			\end{alignat*}
			\normalsize
			Consider the inner integral:
			\footnotesize
			\begin{align}\label{Inner integral simplified}
				& \int_{\tau_{j-1}}^{1}
				\underbrace{\vphantom{\left(\int_{\tau_j}^1\right)}\mathbf{1}}_
				{f'(\tau_{j+1})} \cdot \underbrace{\left(
					\int_{\tau_{j+1}}^{1} \ldots \int_{\tau_{N}}^{1}
					\int_{\tau_{j-1}}^{\tau_{j+1}} \ldots \cdot A_j
					e^{-(\tau_{j+2} - \tau_{j})H} A_{j+1} \cdot \ldots \,
					\dint
					\tau_j \, \dint \tau_{N+1} \, \ldots
					\dint \tau_{j+2} \right)}_{g(\tau_{j+1})} \, \dint
				\tau_{j+1}
				\notag \\
				& \qquad = f(\tau_{j+1})g(\tau_{j+1})\Big \vert_{\tau_{j+1} =
					\tau_{j-1}}^{\tau_{j+1} = 1} - \int_{\tau_{j-1}}^1
				f(\tau_{j+1}) g'(\tau_{j+1}) \, \dint \tau_{j+1} \notag \\
				& \qquad = 0 - \int_{\tau_{j-1}}^1 \tau_{j+1} g'(\tau_{j+1})
				\,
				\dint \tau_{j+1}.
			\end{align}
			\normalsize
			We note that $g$ is of the form
			\begin{equation*}
				g(\tau_{j+1}) = \int_{\tau_{j+1}}^1 \tilde{g}(\tau_{j+1},
				\tau_{j+2}) \, \dint
				\tau_{j+2}
			\end{equation*}
			so that
			\begin{equation*}
				g'(\tau_{j+1}) = -\tilde{g}(\tau_{j+1}, \tau_{j+1}) +
				\int_{\tau_{j+1}}^1 \frac{\partial}{\partial x_1}
				\tilde{g}(\tau_{j+1}, \tau_{j+2}) \, \dint \tau_{j+2}.
			\end{equation*}
			The first summand equals
			\footnotesize
			\begin{equation*}
				-\int_{\tau_{j+1}}^1 \ldots \int_{\tau_{N}}^1
				\int_{\tau_{j-1}}^{\tau_{j+1}} \ldots \cdot A_j
				e^{-(\tau_{j+1} - \tau_{j})H} A_{j+1} e^{-(\tau_{j+3} -
					\tau_{j+1})H} \cdot \ldots \, \dint
				\tau_j \, \dint \tau_{N+1} \, \ldots
				\dint \tau_{j+3}
			\end{equation*}
			\normalsize
			whereas the second equals
			\footnotesize
			\begin{equation*}
				\int_{\tau_{j+1}}^1 \int_{\tau_{j+2}}^1 \ldots
				\int_{\tau_{N}}^1 \ldots \cdot A_{j-1}e^{-(\tau_{j+1} -
					\tau_{j-1})H} A_j e^{-(\tau_{j+2} - \tau_{j+1})H}A_{j+1}
				\cdot
				\ldots
				\, \dint \tau_{N+1} \, \ldots \, \dint \tau_{j+2}.
			\end{equation*}
			\normalsize
			If we plug this into \cref{Inner integral simplified}, we get
			\footnotesize
			\begin{align*}
				\ldots & = \int_{\tau_{j-1}}^1 \tau_{j+1}\int_{\tau_{j+1}}^1
				\ldots
				\int_{\tau_{N}}^1
				\int_{\tau_{j-1}}^{\tau_{j+1}} \ldots \cdot A_{j+1}
				e^{-(\tau_{j+3} -
					\tau_{j+1})H} \cdot \ldots \, \dint
				\tau_j \, \dint \tau_{N+1} \, \ldots
				\dint \tau_{j+3} \, \dint \tau_{j+1} \\
				& \qquad - \int_{\tau_{j-1}}^1 \tau_{j+1} \int_{\tau_{j+1}}^1
				\int_{\tau_{j+2}}^1 \ldots
				\int_{\tau_{N}}^1 \ldots \cdot A_{j-1}e^{-(\tau_{j+1} -
					\tau_{j-1})H} \cdot \ldots
				\, \dint \tau_{N+1} \, \ldots \, \dint \tau_{j+2} \, \dint
				\tau_{j+1} \\
				& = \int_{\tau_{j-1}}^1 \tau_{j+1}\int_{\tau_{j+1}}^1
				\ldots
				\int_{\tau_{N-1}}^1
				\int_{\tau_{j-1}}^{\tau_{j+1}} \ldots \cdot A_{j+1}
				e^{-(\tau_{j+2} -
					\tau_{j+1})H} \cdot \ldots \, \dint
				\tau_j \, \dint \tau_{N} \, \ldots
				\dint \tau_{j+2} \, \dint \tau_{j+1} \\
				& \qquad - \int_{\tau_{j-1}}^1 \tau_{j} \int_{\tau_{j}}^1
				\int_{\tau_{j+1}}^1 \ldots
				\int_{\tau_{N-1}}^1 \ldots \cdot A_{j-1}e^{-(\tau_{j} -
					\tau_{j-1})H} \cdot \ldots
				\, \dint \tau_{N} \, \ldots \, \dint \tau_{j+1} \, \dint
				\tau_{j} \\
				& = \int_{\tau_{j-1}}^1 \int_{\tau_j}^1
				\tau_{j+1}\int_{\tau_{j+1}}^1
				\ldots \int_{\tau_{N-1}}^1 \ldots \cdot A_{j+1}
				e^{-(\tau_{j+2}
					- \tau_{j+1})H} \cdot \ldots \, \dint \tau_{N} \, \ldots
				\, \dint \tau_{j+1} \, \dint \tau_j \\
				& \qquad - \int_{\tau_{j-1}}^1 \tau_{j} \int_{\tau_{j}}^1
				\int_{\tau_{j+1}}^1 \ldots
				\int_{\tau_{N-1}}^1 \ldots \cdot A_{j-1}e^{-(\tau_{j} -
					\tau_{j-1})H} \cdot \ldots
				\, \dint \tau_{N} \, \ldots \, \dint \tau_{j+1} \, \dint
				\tau_{j}
			\end{align*}
			\normalsize
			Taking the outer integral over this expression, we finally have
			\begin{equation}\label{Pulling out 1}
				\{A_1, \ldots, A_j, \mathbf{1}, A_{j+1}, \ldots, A_N\}_H =
				\int_{\Delta_N} (\tau_{j+1} - \tau_j) e^{-\tau_1 H}
				\prod_{j=1}^{N}
				A_j e^{-(\tau_{j+1} - \tau_j)H} \, \mathrm{d} \tau
			\end{equation}
			as desired.
		\end{proof}
		Since $\sum_{j=0}^{N} (\tau_{j+1} - \tau_j) = 1$, this yields
		\begin{equation*}
			\sum_{j = 0}^N \{A_1, \ldots, A_j, \mathbf{1}, A_{j+1}, \ldots,
			A_N\}_H
			= \{A_1, \ldots, A_N\}_H
		\end{equation*}
		which in combination with \cref{Trace cyclic with 1} proves the
		statement.
	\end{proof}
	Since all operators in question supercommute with the Clifford action,
	we may apply \cref{Clifford trace cyclic sum property} to \cref{Chern
		pullback long sum} which yields
	\begin{equation*}
		\ldots = - \CStr(\mathbf{c}(\theta_0')
		\underline{\Phi}_1^{\mathscr{M}_{\mathbb{T}}}(\theta_1, \ldots,
		\theta_N)) = -
		\mathrm{Ch}_{\mathscr{M}_{\mathbb{T}}}(\theta_0,
		\ldots, \theta_N)
	\end{equation*}
	which proves the lemma.
\end{proof}
By definition of the graded dual of a map and due to \cref{Chern even/odd},
we have
\begin{equation*}
	(\underline{d}_{\mathbb{T}} + \underline{b} - \underline{B})^{\vee}
	\mathrm{Ch}_{\mathscr{M}_{\mathbb{T}}} = (-1)^q
	\mathrm{Ch}_{\mathscr{M}_{\mathbb{T}}} \circ
	(\underline{d}_{\mathbb{T}} + \underline{b} - \underline{B})
\end{equation*}
where $\underline{d}_{\mathbb{T}} \coloneqq \underline{d_{\T}}$ is the
first differential on $\mathrm{C}(\Omega_{\mathbb{T}})$ induced by $d_\T$.
In order to employ \cref{Ch pullback}, we have to calculate the composition
of $\alpha$ with the various differentials.
A standard calculation shows that\footnote{This is precisely \cite[Equation
(5.7)]{GüneysuLudewig} modulo the fact that there is a small
misprint in the equation concerning $\underline{d}_{\mathbb{T}}\alpha$ as
well as an ambiguity in the definition of $\underline{S}$.}
\begin{equation*}
	\underline{b}' \alpha = \alpha \underline{b}, \quad \alpha
	\underline{d}_{\mathbb{T}} = \underline{d}_{\mathbb{T}} \alpha - h,
	\quad \alpha \underline{B} = \underline{S}' h,
\end{equation*}
where in the first equation $\underline{b}'$ denotes the second differential
on the quotient
$\mathrm{B}(\underline{\Omega_{\mathbb{T}}})$, in the second equation on the
right hand side $\underline{d}_{\mathbb{T}}$ denotes the first
differential on $\mathrm{B}(\underline{\Omega_{\mathbb{T}}})$ and $h$ is
given by
\begin{equation*}
	h \colon \mathrm{C}(\Omega_{\mathbb{T}}) \to
	\mathrm{B}(\underline{\Omega_{\mathbb{T}}}), \quad (\theta_0, \ldots,
	\theta_N) \mapsto \underline{\mathbf{N}}(\theta_0, \ldots, \theta_N)
\end{equation*}
and $\underline{S}'$ is given by
\begin{equation*}
	\underline{S}'(\theta_1, \ldots, \theta_N) \coloneqq (\sigma,
	\theta_1,
	\ldots, \theta_N) + \underline{S}(\theta_1, \ldots, \theta_N)
\end{equation*}
where $\underline{S}$ is defined on
$\mathrm{B}(\underline{\Omega_{\mathbb{T}}})$ analogously to $S$ introduced
in \cref{Chen normalization S}.
This yields
\begin{equation}\label{Interaction alpha}
	\alpha(\underline{d}_{\mathbb{T}} + \underline{b} - \underline{B}) =
	(\underline{d}_{\mathbb{T}} + \underline{b}')\alpha - (\underline{S}' +
	\mathbf{1})h.
\end{equation}
and hence, due to \cref{Ch pullback}, we have
\begin{align*}
	(\underline{d}_{\mathbb{T}} + \underline{b} - \underline{B})^{\vee}
	\mathrm{Ch}_{\mathbb{T}} &= (-1)^{q+1} \CStr\big(
	\underline{\Phi}_1^{\mathscr{M}_{\mathbb{T}}} \alpha
	(\underline{d}_{\mathbb{T}} + \underline{b} - \underline{B})\big) \\
	\notag
	& = (-1)^{q+1} \CStr\big(
	\underline{\Phi}_1^{\mathscr{M}_{\mathbb{T}}}
	(\underline{d}_{\mathbb{T}} + \underline{b}')\alpha\big) - (-1)^{q+1}
	\CStr\big(
	\underline{\Phi}_1^{\mathscr{M}_{\mathbb{T}}}(\underline{S}' +
	\mathbf{1})h\big).
\end{align*}
If $c \in \mathrm{B}(\underline{\Omega_{\mathbb{T}}})$, it follows from
\cite[Theorem 4.7]{GüneysuLudewig}
that $\underline{\Phi}_1^{\mathscr{M}_{\mathbb{T}}}
\left((\underline{d}_{\mathbb{T}} + \underline{b}')\alpha(c)\right)$ is a
sum of graded commutators of operators on $\HH$ on which, by virtue of
\cref{Clifford trace is trace}, the Clifford supertrace $\CStr$ vanishes.
Moreover, using that $F(\sigma) = \mathbf{1}$ and $F(\sigma, \theta) = 0
= F(\theta,
\sigma)$ for all $\theta \in \Omega_{\mathbb{T}}$, the same arguments
as used in \cref{Ch pullback} together with \cref{Pulling out 1} show that
\begin{equation}\label{Clifford supertrace and S + id}
	\CStr(
	\underline{\Phi}_1^{\mathscr{M}_{\mathbb{T}}}\underline{S}') =
	-
	\CStr(\underline{\Phi}_1^{\mathscr{M}_{\mathbb{T}}})
\end{equation}
when restricted to the subspace
$\mathrm{B}^{\sharp}(\underline{\Omega_{\mathbb{T}}})$ which is the image of
$\underline{\mathbf{N}}$.
Therefore,
\begin{equation*}
	\CStr\big(
	\underline{\Phi}_1^{\mathscr{M}_{\mathbb{T}}} (\underline{S}' +
	\mathbf{1})h\big) = 0
\end{equation*}
which proves coclosedness.\qed
\subsubsection{Proof of \texorpdfstring{\cref{Thm ChM}}{Theorem 5.10}, Part 3}
We need to prove that $\mathrm{Ch}_{\mathscr{M}_{\mathbb{T}}}$ vanishes on
the subcomplex $\mathrm{D}^{\mathbb{T}}(\Omega)$.
First of all, we notice that $\mathbf{c}(\sigma \theta_0) = 0$ so
that $\mathrm{Ch}_{\mathscr{M}_{\mathbb{T}}}$ vanishes on the image of $R$.
Moreover, since $\alpha(S + \mathbf{1}) = (S' + \mathbf{1})
\alpha$, it follows from \cref{Ch pullback} and \cref{Clifford supertrace
	and S + id} that $\mathrm{Ch}_{\mathscr{M}_{\mathbb{T}}}$ vanishes on the
image of $S + \mathbf{1}$.
Note that due to \cref{Multiplicativity}, for $f \in \Omega^0$ we have
\begin{equation*}
	F_{\mathscr{M}_\T}(f) = 0 \quad \text{and} \quad F_{\mathscr{M}_\T}(f,
	\theta) = F_{\mathscr{M}_\T}(\theta, f) = 0
\end{equation*}
for all $\theta \in \Omega_{\mathbb{T}}$, hence $\ChernT(S_i^{(f)}) = 0$.
Finally, we have to show that $\ChernT$ vanishes on the image of
$T_i^{(f)}$ for all $f \in \Omega^0$ and $i \in \N_0$.
However, this can be proven analogously to \cite[Theorem
5.6]{GüneysuLudewig}.\qed
\subsubsection{Proof of \texorpdfstring{\cref{Thm ChM}}{Theorem 5.10}, Part
4}
We will adapt the proof of \cite[Theorem 6.2]{GüneysuLudewig} to
our more general setting and highlight only the main differences.
To this end, let $\mathscr{M}^s = (\HH, \mathbf{c}^s, Q_s)$, $s\in [0,1]$, be
a homotopy of $\vartheta$-summable $\mathcal{C}_q$-Fredholm modules.
We define a family of
$\Lopclifqtrace(\HH)$-valued bar
cochains:
\begin{equation*}
	\Psi_T^{\mathscr{M}^s} \coloneqq -T\int_0^1
	{\Phi}_{uT}^{\mathscr{M}^s}\, \dot{\omega}_{\mathscr{M}^s}
	{\Phi}_{(1-u)T}^{\mathscr{M}^s} \,\mathrm{d} u.
\end{equation*}
For improved readability, we will write $\Psi^s_T \coloneqq
\Psi_T^{\mathscr{M}^s}$, $\Phi_T^s \coloneqq
\Phi_T^{\mathscr{M}^s}$ etc.
\begin{lem}
	There exists a continuous seminorm $\nu$ on $\Omega$ such that for each
	$s \in [0,1]$ and $T > 0$ and all $\theta_1, \ldots, \theta_N \in \Omega$
	the operator $\Psi_T^{\mathscr{M}^s}(\theta_1, \ldots, \theta_N)$ is
	well-defined and trace-class with
	\begin{equation}\label{Estimate Psi}
		\lVert \Psi_T^s(\theta_1, \ldots, \theta_N) \rVert_1 \leq e^{T/2}
		\Tr(e^{-TQ_s^2/2}) \dfrac{T^{(N+1)/2}}{\lfloor N/2 \rfloor !}
		\nu(\theta_1) \cdot
		\ldots \cdot \nu(\theta_N).
	\end{equation}
	The same holds true for $\Psi_T^sQ_s$ and $Q_s \Psi_T^s$ instead
	of $\Psi_T^s$ if we replace
	$T^{(N+1)/2}$ with $T^{N/2}$ and $\lfloor N/2\rfloor!$ with $\lfloor
	(N-1)/2 \rfloor!$ on the right-hand side.
\end{lem}
\begin{proof}
	Similarly to \cref{Fundamental estimate}, it suffices to prove this for
	$\Psi_T^s$.
	We need the following
	\begin{claim}
		For suitable operators $A_1, \ldots, A_N, B$ on $\HH$, it holds that
		\begin{align}\label{Bracket B}
			\begin{split}
				&\{A_1, \ldots, A_{k-1}, B, A_k, \ldots, A_N\}_{TQ_s^2} \\
				& ~~~~~~~~~~= \int_0^1 u^{k-1}(1-u)^{N-(k-1)}\{A_1, \ldots,
				A_{k-1}\}_{uTQ_s^2} B \{A_k, \ldots, A_N\}_{(1-u)TQ_s^2} \,
				\mathrm{d} u.
			\end{split}
		\end{align}
	\end{claim}
	\begin{proof}[Proof of Claim]
		For an integrable function $G \colon \Delta_{N+1} \to E$ in a Banach
		space $E$, we may write
		\footnotesize
		\begin{align*}
			& \int_{\Delta_{N+1}} G(\tau_1, \ldots, \tau_{N+1}) \, \mathrm{d}
			\tau
			\\
			& \qquad = \int_0^1 \dint \tau_{N+1} \int_0^{\tau_{N+1}} \dint
			\tau_N \ldots \int_0^{\tau_2} \dint \tau_1 \, G(\tau_1, \ldots,
			\tau_{N+1}) \\
			& \qquad = \int_0^1 \dint \tau_k \int_{\tau_k}^1 \dint \tau_{N+1}
			\int_{\tau_k}^{\tau_{N+1}} \dint \tau_N \ldots
			\int_{\tau_k}^{\tau_{k+2}} \dint
			\tau_{k+1} \int_0^{\tau_k} \dint \tau_{k-1} \ldots \int_0^{\tau_2}
			\dint \tau_1 \, G(\tau_1, \ldots, \tau_{N+1}) \\
			& \qquad = \int_0^1 \dint t \int_{0}^{1-t} \dint s_{N+1 - k}
			\int_{0}^{s_{N+1 -k}} \dint s_{N-k} \ldots
			\int_{0}^{s_2} \dint s_1 \\
			& ~~~~~~~~~~~~~~~~~~~~~~~~~~~~~~~~~~~~~ \int_0^{t} \dint
			\tau_{k-1} \ldots \int_0^{\tau_2}
			\dint \tau_1 \, G(\tau_1, \ldots, \tau_{k-1}, t, t+s_1, \ldots, t
			+ s_{N+1 - k}).
		\end{align*}
		\normalsize
		Setting
		\footnotesize
		\begin{equation*}
			G(\tau_1, \ldots, \tau_{N+1}) \coloneqq e^{-\tau_1 TQ^2}A_1 \cdot
			\ldots \cdot e^{-(\tau_k - \tau_{k-1})TQ^2} B
			e^{-(\tau_{k+1} - \tau_k)TQ^2} \cdot \ldots \cdot A_N e^{-(1 -
				\tau_{N+1})TQ^2},
		\end{equation*}
		\normalsize
		we see that
		\begin{align*}
			& G(\tau_1, \ldots, \tau_{k-1}, t, t+s_1, \ldots, t +
			s_{N+1 - k}) \\
			& \qquad = \left(e^{-\tau_1 TQ^2}A_1
			e^{-(\tau_2 - \tau_1)TQ^2} \cdot \ldots \cdot A_{k-1} e^{-(t -
				\tau_{k-1})TQ^2}\right) \cdot B \cdot \\
			& \qquad \qquad \cdot \left(e^{-s_1 TQ^2} A_k e^{-(s_2 - s_1)TQ^2}
			\cdot \ldots \cdot A_N e^{-(1 - t - s_{N+1 - k})TQ^2}\right).
		\end{align*}
		Now, with the substitution $\tilde{s}_j = s_j/(1-t)$, the outer
		integral becomes
		\footnotesize
		\begin{align*}
			& \int_{0}^{1-t} \dint s_{N+1 - k}
			\int_{0}^{s_{N+1 -k}} \dint s_{N-k} \ldots
			\int_{0}^{s_2} \dint s_1 \, e^{-s_1 TQ^2} A_k e^{-(s_2 - s_1)TQ^2}
			\cdot \ldots \cdot A_N e^{-(1 - t - s_{N+1 - k})TQ^2} \\
			& \qquad = \int_{0}^{1-t} \dint s_{N+1 - k}
			\int_{0}^{s_{N+1 -k}} \dint s_{N-k} \ldots \\
			& ~~~~~~~~~~~~~~~~~~~~~~~~~ \ldots
			\int_{0}^{s_2} \dint s_1 \, e^{-\frac{s_1}{1-t}(1-t)TQ^2} A_k
			e^{-(\frac{s_2}{1-t} - \frac{s_1}{1-t})(1-t)TQ^2}
			\cdot \ldots \cdot A_N e^{-(1 - \frac{s_{N+1 -
						k}}{1-t})(1-t)TQ^2} \\
			& \qquad = (1-t)^{N+1-k} \int_0^1 \dint \tilde{s}_{N+1-k}
			\int_0^{\tilde{s}_{N+1-k}} \dint \tilde{s}_{N-k} \ldots \\
			& ~~~~~~~~~~~~~~~~~~~~~~~~~ \ldots
			\int_0^{\tilde{s}_2} \dint \tilde{s}_1 \,
			e^{-\tilde{s_1}(1-t)TQ^2}A_k e^{-(\tilde{s_2} -
				\tilde{s_1})(1-t)TQ^2} \cdot \ldots \cdot A_N e^{-(1 -
				\tilde{s}_{N+1-k})(1-t)TQ^2} \\
			& \qquad = (1-t)^{N+1 - k} \{A_k, \ldots, A_N\}_{(1-t)TQ^2}.
		\end{align*}
		\normalsize
		Moreover, with the substitution $\tilde{\tau}_j = \tau_j/t$ the inner
		integral becomes
		\footnotesize
		\begin{align*}
			& \int_0^{t} \dint \tau_{k-1}
			\ldots \int_0^{\tau_2}
			\dint \tau_1 \, e^{-\tau_1 TQ^2}A_1
			e^{-(\tau_2 - \tau_1)TQ^2} \cdot \ldots \cdot A_{k-1} e^{-(t -
				\tau_{k-1})TQ^2} \\
			& \qquad = \int_0^{t} \dint \tau_{k-1}
			\ldots \int_0^{\tau_2}
			\dint \tau_1 \, e^{-\frac{\tau_1}{t} tTQ^2}A_1
			e^{-(\frac{\tau_2}{t} - \frac{\tau_1}{t})tTQ^2} \cdot \ldots
			\cdot A_{k-1} e^{-(1 - \frac{\tau_{k-1}}{t})tTQ^2} \\
			& \qquad = t^{k-1} \int_0^1 \dint \tilde{\tau}_{k-1}
			\int_0^{\tilde{\tau}_{k-1}} \dint \tilde{\tau}_{k-2} \ldots
			\int_0^{\tilde{\tau}_2} \dint \tilde{\tau}_1 \,
			e^{-\tilde{\tau}_1 tTQ^2} A_1 e^{-(\tilde{\tau}_2 -
				\tilde{\tau}_1)tTQ^2} \cdot \ldots \cdot A_{k-1} e^{-(1 -
				\tilde{\tau}_{k-1})tTQ^2} \\
			& \qquad = t^{k-1} \{A_1, \ldots, A_{k-1}\}_{tTQ^2}.
		\end{align*}
		\normalsize
		Finally, plugging those two results into the integral over
		$\Delta_{N+1}$ finishes the proof.
	\end{proof}
	Due to the analytic requirements (H2) and (H3) from \cref{Definition
		homotopy}, the operators  $\dot{\mathbf{c}}^s(\theta)$ and $\dot{Q}_s$
	satisfy the assumptions from \cite[Lemma 4.2]{GüneysuLudewig} so that
	$\dot{\omega}_{\mathscr{M}^s}\Phi_{(1-u)T}^{\mathscr{M}^s}$, and
	consequently the integrand in the definition of $\Psi^s_T$, is a
	well-defined $\Lopclifqtrace(\HH)$-valued bar cochain.
	Applying \cref{Bracket B} together with \cref{Product bar cochains,Defn
		Quantization}, it follows that $\Psi_T^s$ is well-defined
	and $\Psi_T^s(\theta_1, \ldots, \theta_N)$ is equal to
	\footnotesize
	\begin{align*}
		& \sum_{M = 1}^{N} \Bigg((-T)^{M+1} \sum_{j=1}^{M+1} \sum_{I \in
			\mathcal{P}_{M, N}} \{F_s^{\geq 1}(\theta_{I_1}),
		\ldots, F_s^{\geq 1}(\theta_{I_{j-1}}), \dot{Q}_s, F_s^{\geq
			1}(\theta_{I_j}), \ldots, F_s^{\geq 1}(\theta_{I_M})\}_{TQ_s^2} \\
		& \quad + (-T)^M\sum_{j=1}^{M} \sum_{k=1}^N \sum_{\substack{I \in
				\mathcal{P}_{M, N} \\
				I_j =  \{k\}}} \{F_s^{\geq 1}(\theta_{I_1}),
		\ldots, F_s^{\geq 1}(\theta_{I_{j-1}}), \dot{\mathbf{c}}^s(\theta_k),
		F_s^{\geq 1}(\theta_{I_{j+1}}), \ldots, F_s^{\geq
			1}(\theta_{I_M})\}_{TQ_s^2} \Bigg).
	\end{align*}
	\normalsize
	Using the uniform norm bound on $\dot{Q}_s \Delta_s^{-1/2}$, we may apply
	\cite[Lemma 4.2]{GüneysuLudewig} and repeat the steps from the proof of
	\cite[Theorem 4.6]{GüneysuLudewig} to estimate the two brackets in the
	above display.
	The inner sums over give extra factors $N$ and $N^2$ respectively, which
	may be absorbed into the seminorm since $N^2 + N \leq 4^N$ whereas the
	factor $T^{(N+1)/2}$ is a result of the additional $T$ appearing in the
	definition of $\Psi^s_T$.
	Since $\nu$ may be chosen independently of $s$ due to the analytic
	requirements (H1) and (H3) from \cref{Definition homotopy}, we arrive at
	\cref{Estimate Psi}.
\end{proof}
Finally, denote by $\mathscr{M}^s_\T \coloneqq (\HH,\mathbf{c}_{\mathbb{T}}^s,
Q_s), s \in [0,1]$, the homotopy of $\vartheta$-summable
\textit{weak} $\mathcal{C}_q$-Fredholm modules over $\Omega_{\mathbb{T}}$
associated to the homotopy $\mathscr{M}^s$.
The \textit{Chern-Simons form} is then given by
\begin{equation*}
	\mathrm{CS}\big((\mathscr{M}_{\mathbb{T}}^s)_{s \in [0,1]}\big) \coloneqq
	(-1)^q\int_0^1 \alpha^* \CStr\big(
	\underline{\Psi}_1^{\mathscr{M}^s_{\mathbb{T}}}\big) \, \mathrm{d}s
\end{equation*}
where $\alpha$ was defined in \cref{Definition alpha} and
$\underline{\Psi}_1^{\mathscr{M}^s_{\mathbb{T}}}$ is the quotient map of
$\Psi_1^{\mathscr{M}^s_{\mathbb{T}}}$ on
$\mathrm{B}(\underline{\Omega_{\mathbb{T}}})$.
\cref{Estimate Psi} and (H1) from \cref{Definition homotopy} show that this
is a well-defined analytic cochain.
Moreover, since $\Psi_1^{\mathscr{M}^s_{\mathbb{T}}}$ is odd, $\alpha$ is even
and $\CStr$ is even for $q$ even and odd for $q$ odd, it follows that the
Chern-Simons form is odd for $q$ even and even for $q$ odd.
Calculations similar to the ones for $\ChernT$ show that it vanishes on the
images of the operators in \cref{Chen operators} since the $\mathscr{M}^s$
satisfy \cref{Multiplicativity}.
Now, homotopy invariance is a consequence of
\begin{thm}
	The following transgression formula holds:
	\begin{equation*}
		\mathrm{Ch}_{\mathscr{M}^1_{\mathbb{T}}} -
		\mathrm{Ch}_{\mathscr{M}^0_{\mathbb{T}}} = (\underline{d}_{\mathbb{T}}
		+ \underline{b} -
		\underline{B})^{\vee}\mathrm{CS}\big((\mathscr{M}_{\mathbb{T}}^s)_{s
			\in [0,1]}\big).
	\end{equation*}
\end{thm}
\begin{rem}
	Since the $\mathscr{M}^s$ satisfy \cref{Multiplicativity}, it follows
	that $\mathrm{Ch}_{\mathscr{M}^s_{\mathbb{T}}}$ is Chen normalized for
	all $s \in [0,1]$ due to Part 3 of \cref{Thm ChM}.
	Hence, the above formula shows a posteriori that
	$\mathrm{CS}\big((\mathscr{M}_{\mathbb{T}}^s)_{s \in [0,1]}\big)$
	vanishes on all of $\mathrm{D}^{\mathbb{T}}(\Omega)$ and is therefore
	Chen normalized as well.
\end{rem}
\begin{proof}
	According to \cref{Ch pullback}, we have
	\begin{equation*}
		\mathrm{Ch}_{\mathscr{M}^s_{\mathbb{T}}} = -\alpha^*
		\CStr(\underline{\Phi}_1^{\mathscr{M}^s_{\mathbb{T}}}).
	\end{equation*}
	As in \cite[Proposition 6.4]{GüneysuLudewig}, one has
	\begin{equation}\label{Bianchi psi}
		\dfrac{d}{ds} \Phi_T^s = \delta \Psi_T^s + [\omega_s, \Psi_T^s].
	\end{equation}
	Thus,
	\begin{align*}
		\frac{d}{ds}\mathrm{Ch}_{\mathscr{M}^s_{\mathbb{T}}} & = -
		\CStr\Big(\frac{d}{ds}
		\underline{\Phi}_1^{\mathscr{M}_{\mathbb{T}}^s} \alpha\Big) = -
		\left(\CStr\big(\delta \underline{\Psi}_1^{\mathscr{M}_{\mathbb{T}}^s}
		\alpha\big) + \CStr\big([\omega^s,
		\underline{\Psi}_1^{\mathscr{M}_{\mathbb{T}}^s}]\alpha\big)\right) \\
		& = - \CStr
		\big(\underline{\Psi}_1^{\mathscr{M}_{\mathbb{T}}^s}(\underline{d}_\T
		+
		\underline{b}')\alpha\big),
	\end{align*}
	where for the last equality we used that the image of $\alpha$ is
	contained in
	$\mathrm{B}^{\sharp}(\underline{\Omega_{\mathbb{T}}})$ so that $[\omega^s,
	\underline{\Psi}_1^{\mathscr{M}_{\mathbb{T}}^s}]\alpha$
	is a sum of supercommutators on which $\CStr$ vanishes.
	Therefore, with a view on \cref{Interaction alpha}, we get
	\begin{equation*}
		\frac{d}{ds}\mathrm{Ch}_{\mathscr{M}^s_{\mathbb{T}}} = (-1)^{q}
		(\underline{d}_{\mathbb{T}} + \underline{b} - \underline{B})^{\vee}
		\alpha^*\CStr\big(\underline{\Psi}_1^{\mathscr{M}_{\mathbb{T}}^s}\big)
		- \CStr\big(\underline{\Psi}_1^{\mathscr{M}_{\mathbb{T}}^s}
		(\underline{S}' + \mathbf{1})h\big).
	\end{equation*}
	Note that the second summand is zero similarly to \cref{Clifford 
	supertrace and S + id}.
	Hence, integrating both sides of the above equation yields the desired
	transgression formula.
\end{proof}
\subsection{Comparison to Getzler's odd Chern character}
Let $\mathscr{M} = (\HH, \mathbf{c}, Q)$ be an odd $\vartheta$-summable weak
Fredholm module over a locally convex dg algebra $\Omega$.
We define its Chern character via
\begin{equation*}
	\Chern \coloneqq \mathrm{Ch}_{\tilde{\mathscr{M}}}
\end{equation*}
where $\tilde{\mathscr{M}}$ is the associated $\vartheta$-summable weak
$\mathcal{C}_1$-Fredholm module defined at the end of \cref{Section: odd FM}.
If we assume $\Omega$ to be trivially graded with trivial differential $d =
0$, we will see that the cochain $\Chern$ essentially coincides
with the odd Chern character as defined in \cite[Definition 4.1]{GetzlerOdd}.

Note that in this case we have $\Omega^0 = \Omega$ so
that $\mathbf{c}$ is a representation, i. e. it is multiplicative everywhere.
In particular, the same holds true for $\tilde{\mathbf{c}}$ so that
\begin{equation*}
	F_{\tilde{\mathscr{M}}}(\theta_1, \theta_2) = (-1)^{\lvert
		\theta_1 \rvert}(\tilde{\mathbf{c}}(\theta_1 \theta_2) -
	\tilde{\mathbf{c}}(\theta_1)\tilde{\mathbf{c}}(\theta_2)) = 0
\end{equation*}
for all $\theta_1, \theta_2 \in \Omega$.
Also,
\begin{equation*}
	F_{\tilde{\mathscr{M}}}(\theta_1) = [\tilde{Q},
	\tilde{\mathbf{c}}(\theta_1)] - \tilde{\mathbf{c}}(d \theta_1)
	= [\tilde{Q},
	\tilde{\mathbf{c}}(\theta_1)]
	=
	\begin{pmatrix}
		0 & [Q, \mathbf{c}(\theta_1)]\\
		[Q, \mathbf{c}(\theta_1)] & 0
	\end{pmatrix}
\end{equation*}
for $\theta_1 \in \Omega$ and $T > 0$.
Consequently,
\begin{align*}
	\Phi_T^{\tilde{\mathscr{M}}}(\theta_1, \ldots, \theta_N) & = (-T)^N
	\{F_{\tilde{\mathscr{M}}}(\theta_1), \ldots,
	F_{\tilde{\mathscr{M}}}(\theta_N)\}_{T\tilde{Q}^2} \\
	& = (-T)^N \{[\tilde{Q}, \tilde{\mathbf{c}}(\theta_1)], \ldots,
	[\tilde{Q}, \tilde{\mathbf{c}}(\theta_N)]\}_{T\tilde{Q}^2} \\
	& = (-T)^N \int_{\Delta_N}e^{-\tau_1 T\tilde{Q}^2}[\tilde{Q},
	\tilde{\mathbf{c}}(\theta_1)]e^{-(\tau_2 - \tau_1)T\tilde{Q}^2} \cdot
	\ldots \\
	& \qquad \qquad \qquad \ldots \cdot
	e^{-(\tau_{N} - \tau_{N-1})T\tilde{Q}^2}[\tilde{Q},
	\tilde{\mathbf{c}}(\theta_N)]e^{-(1 - \tau_N)T\tilde{Q}^2} \, \mathrm{d}
	\tau
\end{align*}
for $\theta_1, \ldots, \theta_N \in \Omega$.
\begin{lem}
	For $s_1, \ldots, s_{N+1} \geq 0$ and $\theta_1, \ldots, \theta_N \in
	\Omega$
	we have
	\begin{align*}
		& e^{-s_1\tilde{Q}^2}[\tilde{Q},
		\tilde{\mathbf{c}}(\theta_1)]e^{-s_2\tilde{Q}^2} \cdot
		\ldots \cdot
		e^{-s_N\tilde{Q}^2}[\tilde{Q},
		\tilde{\mathbf{c}}(\theta_N)]e^{-s_{N+1}\tilde{Q}^2} = \\
		& \left(
		\begin{smallmatrix}
			0 & e^{-s_1 Q^2}[Q,
			\mathbf{c}(\theta_1)]e^{-s_2 Q^2} \cdot
			\ldots \cdot
			e^{-s_N Q^2}[Q,
			\mathbf{c}(\theta_N)]e^{-s_{N+1}Q^2} \\
			e^{-s_1 Q^2}[Q,
			\mathbf{c}(\theta_1)]e^{-s_2 Q^2} \cdot
			\ldots \cdot
			e^{-s_N Q^2}[Q,
			\mathbf{c}(\theta_N)]e^{-s_{N+1}Q^2} & 0
		\end{smallmatrix}\right)
	\end{align*}
	for $N \geq 1$ odd and
	\begin{align*}
		& e^{-s_1\tilde{Q}^2}[\tilde{Q},
		\tilde{\mathbf{c}}(\theta_1)]e^{-s_2\tilde{Q}^2} \cdot
		\ldots \cdot
		e^{-s_N\tilde{Q}^2}[\tilde{Q},
		\tilde{\mathbf{c}}(\theta_N)]e^{-s_{N+1}\tilde{Q}^2} = \\
		& \left(
		\begin{smallmatrix}
			e^{-s_1 Q^2}[Q,
			\mathbf{c}(\theta_1)]e^{-s_2 Q^2} \cdot
			\ldots \cdot
			e^{-s_N Q^2}[Q,
			\mathbf{c}(\theta_N)]e^{-s_{N+1}Q^2} & 0 \\
			0 & e^{-s_1 Q^2}[Q,
			\mathbf{c}(\theta_1)]e^{-s_2 Q^2} \cdot
			\ldots \cdot
			e^{-s_N Q^2}[Q,
			\mathbf{c}(\theta_N)]e^{-s_{N+1}Q^2}
		\end{smallmatrix}\right)
	\end{align*}
	for $N \geq 1$ even.
\end{lem}
\begin{proof}
	For $1 \leq j \leq N+1$ we have
	\begin{equation*}
		e^{-s_j \tilde{Q}^2} =
		\begin{pmatrix}
			e^{-s_j Q^2} & 0 \\
			0 & e^{-s_j Q^2}
		\end{pmatrix}
	\end{equation*}
	so that
	\begin{equation*}
		e^{-s_j \tilde{Q}^2}[\tilde{Q},
		\tilde{\mathbf{c}}(\theta_j)]e^{-s_{j+1} \tilde{Q}^2} =
		\begin{pmatrix}
			0 & e^{-s_j Q^2}[Q, \mathbf{c}(\theta_j)]e^{-s_{j+1}Q^2} \\
			e^{-s_j Q^2}[Q, \mathbf{c}(\theta_j)]e^{-s_{j+1}Q^2} & 0
		\end{pmatrix}.
	\end{equation*}
	The claim is an easy consequence of this calculation.
\end{proof}
Since all elements in $\Omega$ are even, it follows that
\begin{equation*}
	\gamma \tilde{\mathbf{c}}(\theta_0) =
	\begin{pmatrix}
		0 & \mathbf{1} \\
		-\mathbf{1} & 0
	\end{pmatrix}
	\cdot
	\begin{pmatrix}
		\mathbf{c}(\theta_0) & 0 \\
		0 & \mathbf{c}(\theta_0)
	\end{pmatrix}
	=
	\begin{pmatrix}
		0 & \mathbf{c}(\theta_0) \\
		-\mathbf{c}(\theta_0) & 0
	\end{pmatrix}.
\end{equation*}
Thus, putting together the above calculations, it follows that
$\mathrm{Ch}_{\mathscr{M}}(\theta_0, \ldots, \theta_N) = 0$ for $N \geq 0$ even
and
\begin{align*}
	\mathrm{Ch}_{\mathscr{M}}(\theta_0, \ldots, \theta_N) & =
	\CStr(\tilde{\mathbf{c}}(\theta_0)
	\Phi_1^{\tilde{\mathscr{M}}}(\theta_1, \ldots, \theta_N)) \\
	& = \frac 12
	\CStr\big(\gamma \tilde{\mathbf{c}}(\theta_0)
	\Phi_1^{\tilde{\mathscr{M}}}(\theta_1, \ldots, \theta_N)) \\
	& = 2 \cdot \frac 12 \Tr\big(\gamma \mathbf{c}(\theta_0)
	\Phi_1^{\mathscr{M}}(\theta_1, \ldots, \theta_N)) \\
	& = (-1)^N\int_{\Delta_N} \Tr(\mathbf{c}(\theta_0)e^{-\tau_1 Q^2}[Q,
	\mathbf{c}(\theta_1)]e^{-(\tau_2 - \tau_1) Q^2} \cdot \ldots \\
	& \qquad \qquad \qquad \ldots \cdot
	e^{-(\tau_N - \tau_{N-1})Q^2}[Q,
	\mathbf{c}(\theta_N)]e^{-(1 - \tau_N)Q^2}\big) \, \mathrm{d} \tau
\end{align*}
for $N \geq 0$ odd which, modulo signs and constants (which are the result of
differing sign conventions and normalizations of the supertrace), is nothing
else but the formula given in \cite[Definition 4.1, Proposition
4.2]{GetzlerOdd}.
Thus, our $\mathrm{Ch}_{\mathscr{M}}$ is a differential graded extension
of the odd Chern character defined by Getzler.
\section{A noncommutative index theorem for the odd Chern character}
Given an odd $\vartheta$-summable Fredholm module $\mathscr{M} = (\HH, \pi,
\mathrm{D})$ over a Banach $*$-algebra $\mathcal{A}$ in the sense of
\cite{ConnesCompact}, one can show that pairing the corresponding odd Chern
character $\mathrm{Ch}_*(\mathrm{D})$ with the odd Chern character
$\mathrm{Ch}^{*}(g)$ of a unitary matrix $g \in \mathrm{U}_m(\mathcal{A})$
yields the spectral flow of $\mathrm{D}$ and the twisted operator
$g^{-1}\mathrm{D}g$:
\begin{thm}[\protect{\cite[Theorem 4.3]{GetzlerOdd}}]\label{Getzler sf}
	We have
	\begin{equation}\label{Spectral flow pairing}
		\langle \mathrm{Ch}_*(\mathrm{D}), \mathrm{Ch}^{*}(g) \rangle =
		\frac{1}{\sqrt{\pi}} \int_0^1 \Tr(\dot{\mathrm{D}}_u
		e^{-\mathrm{D}_u^2}) \, \mathrm{d}u = \mathrm{sf}(D,
		g^{-1}\mathrm{D}g),
	\end{equation}
	where $\mathrm{D}_u \coloneqq (1-u)\mathrm{D} + ug^{-1}\mathrm{D}g$ and
	$\dot{\mathrm{D}}_u \coloneqq g^{-1} [\mathrm{D}, g]$.
\end{thm}
Consider now a locally convex dg algebra $\Omega$, $g \in
\mathrm{GL}_m(\Omega^0)$ and an odd $\vartheta$-summable Fredholm
module $\mathscr{M} = (\HH, \mathbf{c}, Q)$ over $\Omega$ as introduced in
\cref{Definition odd FM}.
We want to construct an odd entire chain $\mathrm{Ch}(g)$ which resembles
the odd Chern character of $g$ in the differential graded setting.
Moreover, we would like to obtain an integral formula similar to
\cref{Getzler sf} for the pairing of the analytic cochain
$\mathrm{Ch}_{\tilde{\mathscr{M}}}$,
where $\tilde{\mathscr{M}}$ is the $\vartheta$-summable
$\mathcal{C}_1$-Fredholm module associated to $\mathscr{M}$, with
$\mathrm{Ch}(g)$.
\subsection{Definition of
	\texorpdfstring{$\mathrm{Ch}(g)$}{Ch(g)}}\label{Section: definition cherng}
We will follow the constructions given in \cite[Section 5]{GüneysuCacciatori},
which deals with the special case $\Omega = C^{\infty}(M)$ for some smooth
compact manifold $M$, and make the necessary changes to accommodate our more
general setting.
To this end, we will define an odd entire chain $\mathrm{Ch}(g) \in
\mathrm{C}^{\epsilon}_{-}(\Omega_{\mathbb{T}})$ on the acyclic extension
$\Omega_{\mathbb{T}}$ of $\Omega$.\footnote{Note that we follow the sign
	conventions from \cite[Eq. (2.7)]{GüneysuLudewig} for the definition of
	$\Omega_{\mathbb{T}}$, i. e. we
	consider $d_{\mathbb{T}} = d - \iota$ instead of $d + \iota$, the latter
	convention being used in \cite[Section 3]{GüneysuCacciatori}.}
\begin{defn}
	For $g \in \mathrm{GL}_m(\Omega^0)$ we define the \textit{Maurer-Cartan
		form} $\omega \coloneqq \omega_{g} \coloneqq g^{-1} dg \in
	\mathrm{Mat}_m(\Omega^1) \subseteq
	\mathrm{Mat}_m(\Omega_{\mathbb{T}})$.
\end{defn}
\begin{lem}
	The Maurer-Cartan form satisfies the following equation:
	\begin{equation*}
		d_{\mathbb{T}}\omega + \omega^2 = 0.
	\end{equation*}
\end{lem}
\begin{proof}
	Since $\omega \in \mathrm{Mat}_m(\Omega)$, the action of the
	differential $d_{\mathbb{T}}$ is given by $d$.
	Using that $g^{-1}$ has degree $0$, the graded Leibniz rule implies
	\begin{equation*}
		d \omega = d(g^{-1}dg) = (dg^{-1})dg + g^{-1}d^2 g = (dg^{-1})dg.
	\end{equation*}
	On the other hand,
	\begin{equation}\label{dg^-1}
		0 = d(\mathbf{1}) = d(g^{-1}g) = (dg^{-1})g + g^{-1}dg \iff dg^{-1} =
		-g^{-1}(dg)g^{-1}.
	\end{equation}
	Therefore,
	\begin{equation*}
		d \omega = (dg^{-1})dg = -g^{-1}(dg)g^{-1}dg = -(g^{-1}dg)(g^{-1}dg) =
		-\omega^2
	\end{equation*}
	as desired.
\end{proof}
Following \cite[Eq. (24)]{GüneysuCacciatori}, we define
\begin{equation*}
	\mathcal{A}^s(g) \coloneqq s \omega - s(1-s)\sigma \omega^2, \quad
	\mathcal{B}^s(g) \coloneqq \sigma \omega
\end{equation*}
in $\mathrm{Mat}_m(\Omega_{\mathbb{T}})$ for $s \in I \coloneqq
[0,1]$.
Note that we have slightly different signs in the above definitions since we
use the convention $d_{\mathbb{T}} = d - \iota$.
\begin{defn}
	The \textit{odd Chern character of $g$} is given by the infinite cyclic
	chain
	\begin{equation*}
		\mathrm{Ch}(g) \coloneqq (\mathrm{Ch}_0(g), \mathrm{Ch}_1(g), \ldots)
		\in \prod_{N = 0}^{\infty} \Omega_\T \otimes
		\underline{\Omega_\T}[1]^{\otimes N}
	\end{equation*}
	where $\mathrm{Ch}_0(g) \coloneqq 0$ and
	\begin{equation*}
		\mathrm{Ch}_N(g) \coloneqq \Tr \left(\int_0^1 \mathbf{1} \otimes
		\sum_{k=1}^{N} \mathcal{A}^s(g)^{\otimes(k-1)} \otimes
		\mathcal{B}^s(g)
		\otimes \mathcal{A}^s(g)^{\otimes(N-k)} \, \mathrm{d}s\right)
	\end{equation*}
	for $N > 0$ and $\Tr$ denotes the generalized trace map,
	\begin{align*}\label{Definition generalized trace}
		\begin{split}
			\Tr \colon \mathrm{C}^\epsilon(\mathrm{Mat}_n(\Omega_\T)) &\to
			\mathrm{C}^\epsilon(\Omega_\T), \\
			(\theta_0, \ldots, \theta_N) &\mapsto \sum_{i_0, \ldots, i_N =
				1}^n
			\left((\theta_0)_{i_0}^{i_1}, (\theta_1)_{i_1}^{i_2}, \ldots,
			(\theta_N)_{i_N}^{i_0}\right).
		\end{split}
	\end{align*}
\end{defn}
We now prove:
\begin{thm}
	The following assertions hold:
	\begin{enumerate}
		\item One has $\mathrm{Ch}(g) \in
		\mathrm{C}^{\epsilon}_{-}(\Omega_{\mathbb{T}})$.
		\item One has $(\underline{d}_{\mathbb{T}} + \underline{b} -
		\underline{B})\mathrm{Ch}(g) =
		0$ in the Chen normalized entire complex $\mathrm{N}^{\mathbb{T},
			\epsilon}(\Omega)$.
	\end{enumerate}
\end{thm}
\begin{proof}
	1. Given a continuous seminorm $\nu$ on $\Omega_{\mathbb{T}}$, let
	\begin{equation*}
		C_{\nu} \coloneqq \sup_{s \in I} \, \max\left(\nu(\mathbf{1}),
		\max_{i,j =
			1, \ldots, m} \nu(\mathcal{A}^s(g)_{i}^j), \max_{i,j = 1, \ldots,
			m}
		\nu(\mathcal{B}^s(g)_{i}^j)\right).
	\end{equation*}
	With a view on \cref{Definition entire seminorm}, we see that
	\begin{equation}\label{Seminorm Ch(g)}
		\epsilon_{\nu}(\mathrm{Ch}(g)) \leq mC_{\nu}\sum_{N=0}^{\infty} N
		\dfrac{(m C_{\nu})^N}{\lfloor N/2 \rfloor!}.
	\end{equation}
	Since $\sqrt[N]{\lfloor N/2 \rfloor!} \to \infty$ as $N \to \infty$, it
	follows that
	\begin{equation*}
		\limsup_{N \to \infty}\sqrt[N]{\dfrac{N(mC_{\nu})^N}{\lfloor N/2
				\rfloor!}} = \limsup_{N \to \infty} \sqrt[N]{N} (m C_{\nu})
		\dfrac{1}{\sqrt[N]{\lfloor N/2 \rfloor!}} = 0
	\end{equation*}
	so that the sum in \cref{Seminorm Ch(g)} converges.
	Hence, $\mathrm{Ch}(g)$ is entire.
	It is easy to check that $\mathrm{Ch}(g)$ is odd.

	2. We have to prove that $(\underline{d}_{\mathbb{T}} + \underline{b} -
	\underline{B})\mathrm{Ch}(g) \in
	\overline{\mathrm{D}^{\mathbb{T}}(\Omega)}$.

	First of all, we notice that
	$\underline{B}\mathrm{Ch}(g)$ is a sum of elements of the form
	$S_i^{(\mathbf{1})}$ (see \cref{Si f}) so that $\underline{B}
	\mathrm{Ch}(g) \in \overline{\mathrm{D}^{\mathbb{T}}(\Omega)}$.

	Next, we prove that $(\underline{d}_{\mathbb{T}} +
	\underline{b})\mathrm{Ch}(g) \in
	\overline{\mathrm{D}^{\mathbb{T}}(\Omega)}$.
	To this end, by definition of the differentials, it follows that
	\begin{equation*}
		((\underline{d}_{\mathbb{T}} +
		\underline{b})\mathrm{Ch}(g))_N=((\underline{d}_{\mathbb{T}} +
		\underline{b})
		\mathrm{Ch}_N(g))_N+((\underline{d}_{\mathbb{T}} + \underline{b})
		\mathrm{Ch}_{N+1}(g))_N.
	\end{equation*}
	We compute some of the expressions appearing in the
	differentials.
	To this end, starting with $\underline{d}_{\mathbb{T}}$, let
	\begin{equation*}
		\theta^{(k)} \coloneqq \mathbf{1} \otimes {\mathcal
			A}^{s}(g)^{\otimes (k-1)} \otimes {\mathcal B}^s(g) \otimes
		{\mathcal A}^{s}(g)^{\otimes (N-k)}
	\end{equation*}
	for $1 \leq k \leq N$.
	Using the Maurer-Cartan equation, we get
	\begin{align*}
		d_{\mathbb{T}} \mathcal{A}^s(g) =
		d_{\mathbb{T}}(s\omega_g - s(1-s)\sigma \omega_g^2) & = s
		d\omega_g + s(1-s)\sigma d(\omega_g^2) + s(1-s)\omega_g^2 \\
		& = -s^2\omega_g^2
		+ s(1-s)\sigma d(\omega_g^2) = -s^2 \omega_g^2,
	\end{align*}
	where we used that $d(\omega_g^2) = (d \omega_g)\omega_g - \omega_g (d
	\omega_g) = -\omega_g^3 + \omega_g^3 = 0$ for the last equality.
	Furthermore,
	\begin{equation*}
		d_{\mathbb{T}} \mathcal{B}^s(g) =
		d_{\mathbb{T}}(\sigma \omega_g) = -\sigma d\omega_g
		-\omega_g = -(-\sigma \omega_g^2 + \omega_g).
	\end{equation*}
	By definition of the grading on $\mathrm{Mat}_m(\Omega_{\mathbb{T}})$, we
	have $\lvert \mathcal{A}^s(g) \rvert = 1$
	and $\lvert \mathcal{B}^s(g) \rvert = 0$.
	Thus, for $1 \leq k \leq N$ and $1 \leq i \leq k$, we have
	\begin{equation*}
		m_{i-1}^{(k)} \coloneqq m_{i-1}(\theta^{(k)}) = 0 +
		\sum_{j=1}^{i-1} 1 - (i-1) = 0
	\end{equation*}
	and for $i \geq k+1$ we get
	\begin{equation*}
		m_{i-1}^{(k)} = 0 + \sum_{j=1}^{k-1} 1 + 0 +
		\sum_{j=1}^{i-1-k} 1 - (i-1) = 1 \mod 2.
	\end{equation*}
	On the one hand,
	\begin{align*}
		(\underline{d}_\T + \underline{b})\left(\sum_{k=1}^N
		\theta^{(k)}\right)_N
		& =
		\sum_{k=1}^N \underline{d}_\T(\theta^{(k)}) \\
		& = \sum_{k=1}^N\left((d_{\mathbb{T}} \mathbf{1} \otimes \ldots)
		- \sum_{i = 1}^N (-1)^{m_{i-1}^{(k)}} (\underbrace{\mathbf{1} \otimes
			\ldots \otimes}_i d_{\mathbb{T}} \otimes \ldots)\right) \\
		& = - \mathbf{1} \otimes \sum_{k=1}^N
		\sum_{i=0}^{k-2} \mathcal{A}^s(g)^{\otimes i} \otimes (-s^2
		\omega_g^2)
		\otimes \mathcal{A}^s(g)^{\otimes (k-i-2)} \\
		& \qquad \qquad \qquad \qquad \otimes \sigma \omega_g
		\otimes  {\mathcal A}^s(g)^{\otimes (N-k)} \\
		&\quad+ \mathbf{1} \otimes \sum_{k=1}^N
		\sum_{i=0}^{N-k-1} \mathcal A^s(g)^{\otimes (k-1)} \otimes
		\sigma \omega_g \\
		& \qquad \qquad \qquad \qquad \otimes \mathcal A^s(g)^{\otimes i}
		\otimes (-s^2\omega_g^2 ) \otimes
		{\mathcal A}^s(g)^{\otimes (N-k-i-1)} \\
		&\quad + \mathbf{1} \otimes \sum_{k=1}^N \mathcal
		A^s(g)^{\otimes (k-1)} \otimes (-\sigma\omega_g^2
		+\omega_g) \otimes \mathcal A^s(g)^{\otimes (N-k)}.
	\end{align*}
	Therefore,
	\begin{align*}
		((\underline{d} + \underline{b})\mathrm{Ch}(g))_N
		& = - \Tr \left(\int_0^1 \mathbf{1} \otimes \sum_{k=1}^N
		\sum_{i=0}^{k-2} \mathcal{A}^s(g)^{\otimes i} \otimes (-s^2
		\omega_g^2)
		\otimes \mathcal{A}^s(g)^{\otimes (k-i-2)} \right.\\
		& \left. \vphantom{\int_0^1}\qquad \qquad \qquad \qquad \otimes
		\sigma \omega_g
		\otimes  {\mathcal A}^s(g)^{\otimes (N-k)} \, \mathrm{d}s \right)\\
		&\quad+ \Tr\left(\int_0^1 \mathbf{1} \otimes \sum_{k=1}^N
		\sum_{i=0}^{N-k-1} \mathcal A^s(g)^{\otimes (k-1)} \otimes
		\sigma \omega_g \right.\\
		& \left. \vphantom{\int_0^1}\qquad \qquad \qquad \qquad \otimes
		\mathcal A^s(g)^{\otimes i}
		\otimes (-s^2\omega_g^2 ) \otimes
		{\mathcal A}^s(g)^{\otimes (N-k-i-1)} \, \mathrm{d}s \right)\\
		&\quad + \Tr\left(\int_0^1 \mathbf{1} \otimes \sum_{k=1}^N \mathcal
		A^s(g)^{\otimes (k-1)} \otimes (-\sigma\omega_g^2
		+\omega_g) \otimes \mathcal A^s(g)^{\otimes (N-k)} \, \mathrm{d}s
		\right).
	\end{align*}
	On the other hand,
	\begin{equation*}
		(\underline{d}_\T + \underline{b})\left(\sum_{k=1}^{N+1}
		\theta^{(k)}\right)_N = \sum_{k=1}^{N+1} \underline{b}(\theta^{(k)})
	\end{equation*}
	where $\theta^{(k)} \coloneqq \mathbf{1} \otimes \mathcal{A}^s(g)^{\otimes
		(k-1)} \otimes \mathcal{B}^s(g) \otimes
	\mathcal{A}^s(g)^{\otimes(N+1-k)}$
	for $1 \leq k \leq N + 1$.
	Calculating the expressions appearing in $\underline{b}$, we have
	\begin{align*}
		\mathcal{A}^s(g) \mathcal{A}^s(g) & = (s \omega_g - s(1-s)\sigma
		\omega_g^2)(s \omega_g - s(1-s)\sigma \omega_g^2) \\
		& = s^2 \omega_g^2 -
		s\omega_g s(1-s)\sigma \omega_g^2 - s(1-s)\sigma \omega_g^2 s \omega_g
		+ s^2(1-s)^2(\sigma \omega_g^2)^2 \\
		& = s^2 \omega_g^2 +
		s^2(1-s)\sigma \omega_g^3 - s^2(1-s)\sigma \omega_g^3
		+ s^2(1-s)^2\sigma^2 \omega_g^4 \\
		& = s^2 \omega_g^2
	\end{align*}
	and
	\begin{equation*}
		\mathcal{A}^s(g) \mathcal{B}^s(g) = (s \omega_g - s(1-s)\sigma
		\omega_g^2)\sigma \omega_g = s \omega_g \sigma \omega_g = -s
		\sigma\omega_g^2
	\end{equation*}
	as well as
	\begin{equation*}
		\mathcal{B}^s(g) \mathcal{A}^s(g) = s \sigma \omega_g^2.
	\end{equation*}
	With a view on \cref{Definition bunderline}, we calculate (modulo two)
	\begin{equation*}
		\lvert \theta^{(k)}_{N+1} \rvert - 1 =
		\begin{cases}
			0, & \quad 1 \leq k \leq N, \\
			1, & \quad k = N+1
		\end{cases}
	\end{equation*}
	and
	\begin{equation*}
		m_N(\theta^{(k)}) =
		\begin{cases}
			1, & \quad 1 \leq k \leq N, \\
			0, & \quad k = N + 1.
		\end{cases}
	\end{equation*}
	so that for all $k$ we have $(-1)^{(\lvert \theta^{(k)}_{N+1} \rvert -
		1)m_N} = 1$.
	This yields
	\begin{align*}
		((\underline{d}_\T + \underline{b})\mathrm{Ch}_{N+1}(g))_N & = - \Tr
		\left(\int_0^1 \mathbf{1} \otimes \sum_{k=1}^N
		\sum_{i=0}^{k-2} \mathcal{A}^s(g)^{\otimes i} \otimes (+s^2
		\omega_g^2)
		\otimes \mathcal{A}^s(g)^{\otimes (k-i-2)} \right.\\
		& \left. \vphantom{\int_0^1}\qquad \qquad \qquad \qquad \otimes
		\sigma \omega_g
		\otimes  {\mathcal A}^s(g)^{\otimes (N-k)} \, \mathrm{d}s \right)\\
		&\quad+ \Tr\left(\int_0^1 \mathbf{1} \otimes \sum_{k=1}^N
		\sum_{i=0}^{N-k-1} \mathcal A^s(g)^{\otimes (k-1)} \otimes
		\sigma \omega_g \right.\\
		& \left. \vphantom{\int_0^1}\qquad \qquad \qquad \qquad \otimes
		\mathcal A^s(g)^{\otimes i}
		\otimes (+s^2\omega_g^2 ) \otimes
		{\mathcal A}^s(g)^{\otimes (N-k-i-1)} \, \mathrm{d}s \right)\\
		&\quad + \Tr\left(\int_0^1 \mathbf{1} \otimes \sum_{k=1}^N \mathcal
		A^s(g)^{\otimes (k-1)} \otimes 2s\sigma\omega_g^2 \otimes \mathcal
		A^s(g)^{\otimes (N-k)} \, \mathrm{d}s
		\right).
	\end{align*}
	Since $\frac{d}{ds} \mathcal{A}^s(g) = \omega_g + (2s-1)\sigma
	\omega_g^2$, it follows that
	\begin{align*}
		((\underline{d}_\T + \underline{b})\mathrm{Ch}(g))_N & = \Tr
		\left(\int_0^1 \mathbf{1} \otimes \sum_{k = 1}^N
		\mathcal{A}^s(g)^{\otimes(k-1)} \otimes \left(\frac{d}{ds}
		\mathcal{A}^s(g)\right) \otimes \mathcal{A}^s(g)^{\otimes (N-k)} \,
		\mathrm{d}s\right) \\
		& = \Tr\left(\int_0^1 \frac{d}{ds}(\mathbf{1} \otimes
		\mathcal{A}^s(g)^{\otimes N}) \, \mathrm{d}s\right) \\
		& = \Tr\big(\mathbf{1} \otimes \mathcal{A}^1(g)^{\otimes N}\big) -
		\Tr\big(\mathbf{1} \otimes \mathcal{A}^0(g)^{\otimes N}\big) \\
		& = \Tr(\mathbf{1} \otimes \omega_g^{\otimes N}).
	\end{align*}
	\begin{claim}
		Define\footnote{The last part of the proof in
			\cite{GüneysuCacciatori} implicitly uses that the tensor product
			is
			taken over $\Omega^0$ which is wrong.
			The following correction was communicated to me by Sergio
			Cacciatori.}
		\begin{align*}
			\lambda_2 & \coloneqq \Tr(\mathbf{1} \otimes dg^{-1} \otimes dg)
			\\
			\lambda_N & \coloneqq \Tr(\mathbf{1} \otimes
			\omega_g^{\otimes(N-2)}
			\otimes dg^{-1} \otimes dg), \quad N \geq 3
		\end{align*}
		in $\mathrm{C}(\Omega_{\mathbb{T}})$.
		Then,
		\begin{align*}
			\Tr(\mathbf{1} \otimes \omega_g) - \lambda_2 & \in
			\mathrm{D}^{\mathbb{T}}(\Omega)\\
			\Tr(\mathbf{1} \otimes \omega_g^{\otimes N}) + \lambda_N -
			\lambda_{N+1} & \in
			\mathrm{D}^{\mathbb{T}}(\Omega), \quad N \geq 2.
		\end{align*}
	\end{claim}
	\begin{proof}[Proof of Claim]
		For $N = 1$, we have
		\begin{equation*}
			\mathbf{1} \otimes \omega_g = \mathbf{1} \otimes g^{-1}dg =
			g^{-1} \otimes dg + \mathbf{1} \otimes dg^{-1} \otimes dg
		\end{equation*}
		modulo $\mathrm{D}^{\mathbb{T}}(\mathrm{Mat}_m(\Omega))$.
		However, since $g^{-1} \otimes dg = -T_0^{(g)}(g^{-1})$ (see
		\cref{Section: chen normalization} for the definition of $T_0^{(g)}$)
		it follows that
		\begin{equation*}
			\mathbf{1} \otimes \omega_g - \mathbf{1} \otimes dg^{-1} \otimes
			dg \in \mathrm{D}^{\mathbb{T}}(\mathrm{Mat}_m(\Omega))
		\end{equation*}
		which implies $\Tr(\mathbf{1} \otimes \omega_g) - \lambda_2 \in
		\mathrm{D}^{\mathbb{T}}(\Omega)$.

		For $N \geq 2$, we have
		\begin{align*}
			\mathbf{1} \otimes \omega_g^{\otimes N} & = \mathbf{1} \otimes
			\omega_g^{\otimes (N-2)} \otimes g^{-1}dg \otimes g^{-1} dg \\
			& = \mathbf{1} \otimes \omega_g^{\otimes (N-2)} \otimes
			g^{-1}(dg)g^{-1} \otimes dg + \mathbf{1} \otimes
			\omega_g^{\otimes(N-1)} \otimes dg^{-1} \otimes dg \\
			& = -\mathbf{1} \otimes
			\omega_g^{\otimes (N-2)} \otimes dg^{-1} \otimes dg + \mathbf{1}
			\otimes \omega_g^{\otimes(N-1)} \otimes dg^{-1} \otimes dg
		\end{align*}
		modulo $\mathrm{D}^{\mathbb{T}}(\mathrm{Mat}_m(\Omega))$ where we
		used \cref{dg^-1}.
		Hence,
		\begin{equation*}
			\Tr(\mathbf{1} \otimes \omega_g^{\otimes N}) + \lambda_N -
			\lambda_{N+1} \in
			\mathrm{D}^{\mathbb{T}}(\Omega)
		\end{equation*}
		as desired.
	\end{proof}
	Similarly to \cref{Seminorm Ch(g)}, one sees that the $\lambda_N$ form an
	entire chain so that in particular, the sequence $(\lambda_N)_N$
	converges to $0$ in
	$\mathrm{C}^\epsilon(\Omega_{\mathbb{T}})$ as $N \to \infty$.
	Hence, we may write
	\begin{equation*}
		0 = - \lambda_2 + \sum_{N = 2}^\infty (\lambda_N - \lambda_{N+1})
	\end{equation*}
	in $\mathrm{C}^\epsilon(\Omega_{\mathbb{T}})$ which together with above
	claim implies $(\underline{d}_\T +
	\underline{b})\mathrm{Ch}(g) \in
	\overline{\mathrm{D}^{\mathbb{T}}(\Omega)}$.
\end{proof}
\subsection{The spectral flow of an odd Fredholm module}\label{Section:
	spectral
	flow}
We want to make sense of the right-hand side of \cref{Spectral flow pairing}
in our abstract setting.
To this end, let $\mathscr{M} = (\HH, \mathbf{c}, Q)$ be an odd
$\vartheta$-summable Fredholm module over the locally convex dg algebra
$\Omega$ and $\mathscr{M}_{\mathbb{T}}$ its acyclic extension as discussed in
\cref{Acyclic extension FM}.
As before, this gives rise to an odd
$\vartheta$-summable weak Fredholm
module $\mathscr{M}_{\mathbb{T}}^{(m)} = (\HH^m, \mathbf{c}_{\mathbb{T}, m},
Q_m)$ over the locally convex dg algebra
$\mathrm{Mat}_m(\Omega_{\mathbb{T}})$.
Recall that if no grading is specified, we assume all elements to have
degree $0$ so that in particular, the supercommutator $[\cdot, \cdot]$
reduces to the ordinary
commutator whereas the notation $[[\cdot, \cdot]]$ will always be used for
anticommutators.
Define unbounded operators
\begin{align*}
	Q_{g,s} & \coloneqq (1-s)Q_m + s \cdot \mathbf{c}_{\T, m}(g^{-1}) Q_m
	\mathbf{c}_{\T, m}(g) = Q_m + s \cdot \mathbf{c}_{\T, m}(g^{-1})[Q_m,
	\mathbf{c}_{\T, m}(g)]\\
	\dot{Q}_{g,s} & \coloneqq \mathbf{c}_{\T, m}(g^{-1})[Q_m, \mathbf{c}_{\T,
		m}(g)]
\end{align*}
on $\HH^m$ for $s \in I$.
Since $g, g^{-1} \in \mathrm{Mat}_m(\Omega^0)$ it follows
that $[Q_m, \mathbf{c}_{\T, m}(g)] = \mathbf{c}_{\T, m}(d g)$ and therefore,
\begin{equation*}
	\mathbf{c}_{\T, m}(g^{-1})[Q_m, \mathbf{c}_{\T, m}(g)] = \mathbf{c}_{\T,
		m}(g^{-1}) \mathbf{c}_{\T, m}(dg) =
	\mathbf{c}_{\T, m}(g^{-1}dg) = \mathbf{c}_{\T, m}(\omega).
\end{equation*}
Hence,
\begin{equation*}
	Q_{g,s} = Q_m + s \cdot \mathbf{c}_{\T, m}(\omega), \quad \dot{Q}_{g,s} =
	\mathbf{c}_{\T, m}(\omega)
\end{equation*}
and
\begin{equation*}
	Q_{g,s}^2 = Q_m^2 + s[[Q_m, \mathbf{c}_{\T, m}(\omega)]] + s^2
	\mathbf{c}_{\T, m}(\omega)^2 = Q_m^2 + X_s,
\end{equation*}
where $X_s \coloneqq s[[Q_m, \mathbf{c}_{\T, m}(\omega)]] + s^2
\mathbf{c}_{\T, m}(\omega)^2$.
\begin{lem}
	The operator $Q_{g,s}^2$ generates a strongly continuous semigroup of
	operators on $\HH^m$ with infinitesimal generator $-Q_{g,s}^2$ which is
	given the perturbation series
	\begin{equation}\label{Perturbation series}
		e^{-TQ_{g,s}^2} = e^{-T(Q_m^2 + X_s)} = \sum_{M=0}^{\infty} (-T)^M
		\{\underbrace{X_s, \ldots, X_s}_M\}_{TQ^2}
	\end{equation}
	for $T > 0$.
	Moreover, the operator $\dot{Q}_{g,s} e^{-TQ_{g,s}^2}$ is well-defined
	and trace-class.
\end{lem}
\begin{proof}
	Due to the analytic requirement (A1) in \cref{Definition odd FM}, the
	operators $X_s$ satisfy
	the assumptions from \cite[Lemma 4.2]{GüneysuLudewig} with $a_j = \frac
	12$ for $1 \leq j
	\leq M$ and the same is true for $\dot{Q}_{g,s} =
	\mathbf{c}_{\T,m}(\omega)$
	with $a_0 = 0$ since $\mathbf{c}_{\T, m}$ maps to bounded operators.
	Thus,
	\begin{equation*}
		\dot{Q}_{g,s}\{\underbrace{X_s, \ldots, X_s}_{M}\}_{TQ_m^2}
	\end{equation*}
	is a well-defined trace-class operator for each $T > 0$ and the estimates
	from \cite[Lemma 4.2]{GüneysuLudewig} show that the sum in
	\cref{Perturbation series}
	converges in trace norm and defines a strongly continuous semigroup with
	infinitesimal generator $-Q_{g,s}^2$.\footnote{See also
		\cite[Chapter Nine, §2.1]{Kato}.}
\end{proof}
Thus, the integral
\begin{equation*}
	\int_0^1 \Tr(\dot{Q}_{g,s}e^{-Q_{g,s}^2}) \, \mathrm{d}s
\end{equation*}
is well-defined when using the perturbation series for $e^{-Q_{g,s}^2}$.
Note that if $\mathbf{c}(g) \in \Lop(\HH)$ is unitary, the above integral
equals $\sqrt{\pi} \cdot \mathrm{sf}(Q, g^{-1}Qg)$ as shown in
\cite[Corollary 2.7]{GetzlerOdd}.
\subsection{The pairing \texorpdfstring{$\langle \ChernT,
		\mathrm{Ch}(g) \rangle$}{<ChMT, Ch(g)>}}\label{Section: spectral flow
	pairing 1}
Recall that
\begin{align*}
	\mathrm{Ch}_N(g) & \coloneqq \Tr \left(\int_0^1 \mathbf{1} \otimes
	\sum_{k=1}^{N} \mathcal{A}^s(g)^{\otimes(k-1)} \otimes
	\mathcal{B}^s(g)
	\otimes \mathcal{A}^s(g)^{\otimes(N-k)} \, \mathrm{d}s\right)\\
	& = \int_0^1 \Tr \left(\mathbf{1} \otimes
	\sum_{k=1}^{N} \mathcal{A}^s(g)^{\otimes(k-1)} \otimes
	\mathcal{B}^s(g)
	\otimes \mathcal{A}^s(g)^{\otimes(N-k)}\right)\, \mathrm{d}s \\
	& = \int_0^1 \mathbf{1} \otimes \Tr \left(
	\sum_{k=1}^{N} \mathcal{A}^s(g)^{\otimes(k-1)} \otimes
	\mathcal{B}^s(g)
	\otimes \mathcal{A}^s(g)^{\otimes(N-k)}\right) \, \mathrm{d}s
\end{align*}
for $N \geq 1$.
\cref{Doubling acyclic} shows that
$\tilde{\mathscr{M}}_{\mathbb{T}} = \widetilde{\mathscr{M}_{\mathbb{T}}}$,
and it is easy to check that
\begin{equation*}
	U \colon \widetilde{\HH^m} \to \tilde{\HH}^m, \quad (v^{+}, v^{-})
	\mapsto \big((v^{+}_1, v^{-}_{1}), \ldots, (v^{+}_m, v^{-}_m)\big)
\end{equation*}
induces an isomorphism between $\widetilde{\mathscr{M}^{(m)}_{\mathbb{T}}} =
(\widetilde{\HH^m}, \widetilde{\mathbf{c}_{\T, m}}, \widetilde{Q_m})$ and
$\tilde{\mathscr{M}}_{\mathbb{T}}^{(m)} = (\tilde{\HH}^m,
\tilde{\mathbf{c}}_{\T, m}, \tilde{Q}_m)$ as $\vartheta$-summable weak
$\mathcal{C}_1$-Fredholm modules over $\mathrm{Mat}_m(\Omega_{\mathbb{T}})$.
Since their Chern characters
coincide according to \cref{Isomorphism Chern character}, we will identify
the two Fredholm modules from now on.
Hence, using the identity
\begin{equation}\label{Str clif1 tr}
	\mathrm{Str}_{\mathcal{C}_1, \tilde{\HH}}\left((\tilde{\mathbf{c}}_\T
	\Phi_1^{\tilde{\mathscr{M}}_{\mathbb{T}}})(\Tr(\theta_0, \ldots,
	\theta_N))\right) = \mathrm{Str}_{\mathcal{C}_1,
		\tilde{\HH}^m}\left(\tilde{\mathbf{c}}_{\T, m}(\theta_0)
	\Phi_1^{\tilde{\mathscr{M}}_{\mathbb{T}}^{(m)}}(\theta_1, \ldots,
	\theta_N)\right),
\end{equation}
we have
\begin{align}
	&\langle \mathrm{Ch}_{\mathscr{M}_{\mathbb{T}}},
	\mathrm{Ch}(g) \rangle\notag \\
	& \qquad = \sum_{N=0}^{\infty} \int_0^1
	\mathrm{Str}_{\mathcal{C}_1,
		\tilde{\HH}^m}\left(\tilde{\mathbf{c}}_{\mathbb{T}, m}(\mathbf{1})
	\Phi_1^{\tilde{\mathscr{M}}_{\mathbb{T}}^{(m)}}\left(
	\sum_{k=1}^{N} \mathcal{A}^s(g)^{\otimes(k-1)} \otimes
	\mathcal{B}^s(g)
	\otimes \mathcal{A}^s(g)^{\otimes(N-k)}
	\right)\right) \notag\\
	& \qquad = \sum_{N=0}^{\infty} \int_0^1
	\mathrm{Str}_{\mathcal{C}_1, \tilde{\HH}^m}\left(
	\sum_{k=1}^{N}\Phi_1^{\tilde{\mathscr{M}}_{\mathbb{T}}^{(m)}}\left(
	\mathcal{A}^s(g)^{\otimes(k-1)} \otimes
	\mathcal{B}^s(g)
	\otimes \mathcal{A}^s(g)^{\otimes(N-k)}
	\right)\right) \notag\\
	& \qquad = \sum_{N=0}^{\infty} \int_0^1
	\frac 12 \mathrm{Tr}_{\tilde{\HH}^m}\left(\tilde{\gamma}
	\sum_{k=1}^{N}\Phi_1^{\tilde{\mathscr{M}}_{\mathbb{T}}^{(m)}}\left(
	\mathcal{A}^s(g)^{\otimes(k-1)} \otimes
	\mathcal{B}^s(g)
	\otimes \mathcal{A}^s(g)^{\otimes(N-k)}
	\right)\right)\label{Pairing}
\end{align}
where
\begin{equation*}
	\tilde{\gamma} \coloneqq \begin{pmatrix}
		0 & \mathbf{1}_{\HH} \\
		\mathbf{1}_{\HH} & 0 \\
		& & \ddots \\
		& & & 0 & \mathbf{1}_{\HH} \\
		& & & \mathbf{1}_{\HH} & 0\\
	\end{pmatrix} \in \Lop(\tilde{\HH}^m).
\end{equation*}
To proceed further, we have to calculate the expressions involving
$\Phi_1^{\tilde{\mathscr{M}}_{\mathbb{T}}^{(m)}}$:
\begin{align*}
	& \Phi_1^{\tilde{\mathscr{M}}_{\mathbb{T}}^{(m)}}
	\left(\mathcal{A}^s(g)^{\otimes(k-1)} \otimes
	\mathcal{B}^s(g)
	\otimes \mathcal{A}^s(g)^{\otimes(N-k)}\right) \\
	& \qquad = \sum_{M = 1}^N (-1)^M \sum_{J \in \mathcal{P}_{M,N}}
	\{F_{\tilde{\mathscr{M}}_{\mathbb{T}}^{(m)}}(\theta_{J_1}),
	\ldots,
	F_{\tilde{\mathscr{M}}_{\mathbb{T}}^{(m)}}(\theta_{J_M})\}_
	{\tilde{Q}_m^2}
\end{align*}
where $\theta_1 \otimes \ldots \otimes \theta_N \coloneqq
\mathcal{A}^s(g)^{\otimes(k-1)} \otimes \mathcal{B}^s(g)
\otimes \mathcal{A}^s(g)^{\otimes(N-k)}$.
By definition of the bracket, we have to evaluate the following expressions:
\begin{align*}
	& F_{\tilde{\mathscr{M}}_{\mathbb{T}}^{(m)}}(\mathcal{A}^s(g)), \quad
	F_{\tilde{\mathscr{M}}_{\mathbb{T}}^{(m)}}(\mathcal{A}^s(g),
	\mathcal{B}^s(g)), \quad F_{\tilde{\mathscr{M}}_{\mathbb{T}}^{(m)}}(
	\mathcal{B}^s(g), \mathcal{A}^s(g)), \\
	& F_{\tilde{\mathscr{M}}_{\mathbb{T}}^{(m)}}(\mathcal{B}^s(g)), \quad
	F_{\tilde{\mathscr{M}}_{\mathbb{T}}^{(m)}}(
	\mathcal{A}^s(g), \mathcal{A}^s(g)).
\end{align*}
We calculate:
\begin{align*}
	F_{\tilde{\mathscr{M}}_{\mathbb{T}}^{(m)}}(\mathcal{A}^s(g)) & =
	F_{\tilde{\mathscr{M}}_{\mathbb{T}}^{(m)}}(s\omega - s(1-s)\sigma
	\omega^2) \\
	& = s[\tilde{Q}_m,
	\tilde{\mathbf{c}}_{\T, m}(\omega)] - s\tilde{\mathbf{c}}_{\T, m}(d
	\omega) -
	s(1-s)\tilde{\mathbf{c}}_{\T, m}(\omega^2) \\
	& =  s[\tilde{Q}_m,
	\tilde{\mathbf{c}}_{\T, m}(\omega)] + s\tilde{\mathbf{c}}_{\T,
		m}(\omega^2) -
	s(1-s)\tilde{\mathbf{c}}_{\T, m}(\omega^2) \\
	& = s[\tilde{Q}_m,
	\tilde{\mathbf{c}}_{\T, m}(\omega)] +
	s^2\tilde{\mathbf{c}}_{\T, m}(\omega^2)
\end{align*}
where we used the Maurer-Cartan equation.
Note that
\begin{equation*}
	F_{\tilde{\mathscr{M}}_{\mathbb{T}}^{(m)}}(\mathcal{A}^s(g),
	\mathcal{B}^s(g)) = (-1)^{\lvert \omega \rvert} (\tilde{\mathbf{c}}_{\T,
		m}(s\omega
	\cdot 0) - \tilde{\mathbf{c}}_{\T, m}(s \omega) \tilde{\mathbf{c}}_{\T,
		m}(0)) = 0
\end{equation*}
and the same is true for $F(\mathcal{B}^s(g), \mathcal{A}^s(g))$.
Furthermore,
\begin{align*}
	F_{\tilde{\mathscr{M}}_{\mathbb{T}}^{(m)}}(\mathcal{A}^s(g),
	\mathcal{A}^s(g)) & = (-1)^{\lvert \omega
		\rvert}(\tilde{\mathbf{c}}_{\T, m}(s^2 \omega^2) - s^2
	\tilde{\mathbf{c}}_{\T, m}(\omega)
	\tilde{\mathbf{c}}_{\T, m}(\omega)) \\
	& = -(s^2\tilde{\mathbf{c}}_{\T, m}(\omega^2) - s^2
	\tilde{\mathbf{c}}_{\T, m}(\omega)^2)
\end{align*}
and finally,
\begin{equation*}
	F_{\tilde{\mathscr{M}}_{\mathbb{T}}^{(m)}}(\mathcal{B}^s(g)) =
	F_{\tilde{\mathscr{M}}_{\mathbb{T}}^{(m)}}(\sigma \omega) =
	\tilde{\mathbf{c}}_{\T, m}(\omega).
\end{equation*}
Define $\tilde{X}_s \coloneqq s[\tilde{Q}_m, \tilde{\mathbf{c}}_{\mathbb{T},
	m}(\omega)] + s^2\tilde{\mathbf{c}}_{\mathbb{T}, m}(\omega)^2$.
Using the above results, we notice that
\begin{align*}
	F_{\tilde{\mathscr{M}}_{\mathbb{T}}^{(m)}}(\mathcal{A}^s(g),
	\mathcal{A}^s(g)) +
	F_{\tilde{\mathscr{M}}_{\mathbb{T}}^{(m)}}(\mathcal{A}^s(g)) & = s^2
	\tilde{\mathbf{c}}_{\mathbb{T}, m}(\omega)^2 -
	s^2\tilde{\mathbf{c}}_{\mathbb{T}, m}(\omega^2) \\
	& \qquad + s[\tilde{Q}_m, \tilde{\mathbf{c}}_{\mathbb{T}, m}(\omega)] +
	s^2\tilde{\mathbf{c}}_{\mathbb{T}, m}(\omega^2) \\
	& = s[\tilde{Q}_m,
	\tilde{\mathbf{c}}_{\mathbb{T}, m}(\omega)] + s^2
	\tilde{\mathbf{c}}_{\mathbb{T}, m}(\omega)^2 =
	\tilde{X}_s.
\end{align*}
\begin{lem}
	The following identity holds:
	\begin{equation*}
		\sum_{M = 0}^{\infty} (-T)^M \{\underbrace{\tilde{X}_s, \ldots,
			\tilde{X}_s}_M\}_{T\tilde{Q}^2} = \sum_{N=0}^{\infty}
		\Phi_T^{\tilde{\mathscr{M}}_{\mathbb{T}}^{(m)}}
		(\underbrace{\mathcal{A}^s(g),
			\ldots, \mathcal{A}^s(g)}_N).
	\end{equation*}
\end{lem}
\begin{proof}
	Using multi-linearity of the bracket, we have
	\footnotesize
	\begin{align*}
		& \sum_{M = 0}^\infty (-T)^M \{\underbrace{\tilde{X}_s, \ldots,
			\tilde{X}_s}_M\}_{T\tilde{Q}^2} \\
		& \quad = \sum_{M = 0}^\infty (-T)^M
		\{F_{\tilde{\mathscr{M}}_{\mathbb{T}}^{(m)}}(\mathcal{A}^s(g),
		\mathcal{A}^s(g)) +
		F_{\tilde{\mathscr{M}}_{\mathbb{T}}^{(m)}}(\mathcal{A}^s(g)), \ldots,
		F_{\tilde{\mathscr{M}}_{\mathbb{T}}^{(m)}}(\mathcal{A}^s(g),
		\mathcal{A}^s(g)) +
		F_{\tilde{\mathscr{M}}_{\mathbb{T}}^{(m)}}(\mathcal{A}^s(g))\}
		_{T\tilde{Q}^2} \\
		& \quad = \sum_{M = 0}^\infty \sum_{l=0}^M (-T)^M \sum_{J \in
			\mathcal{P}_M^{(l)}}
		\{F_{\tilde{\mathscr{M}}_{\mathbb{T}}^{(m)}}(\theta_{J_1}),
		\ldots, F_{\tilde{\mathscr{M}}_{\mathbb{T}}^{(m)}}(\theta_{J_M})\}
		_{T\tilde{Q}^2}
	\end{align*}
	\normalsize
	where $\mathcal{P}_M^{(l)}$ is the set of partitions $J \in
	\mathcal{P}_{M, 2(M-l)+l}$ such that there exist $i_1, \ldots, i_l \in
	\{1, \ldots, M\}$ for which $\theta_{J_{r}} = \mathcal{A}^s(g)$ for $r
	\in \{i_1, \ldots, i_l\}$ and $\theta_{J_{r}} = (\mathcal{A}^s(g),
	\mathcal{A}^s(g))$ for $r \in \{1, \ldots, M\} \setminus \{i_1,
	\ldots, i_l\}$.
	Substituting $K = 2(M-l) + l$ yields
	\begin{align*}
		& \sum_{l=0}^M (-T)^M \sum_{J \in
			\mathcal{P}_M^{(l)}}
		\{F_{\tilde{\mathscr{M}}_{\mathbb{T}}^{(m)}}(\theta_{J_1}),
		\ldots, F_{\tilde{\mathscr{M}}_{\mathbb{T}}^{(m)}}(\theta_{J_M})\}
		_{T\tilde{Q}^2} \\
		& \qquad = (-T)^M
		\sum_{K=M}^{2M} \sum_{J \in \mathcal{P}_{M, K}}
		\{F_{\tilde{\mathscr{M}}_{\mathbb{T}}^{(m)}}(\theta_{J_1}),
		\ldots, F_{\tilde{\mathscr{M}}_{\mathbb{T}}^{(m)}}(\theta_{J_M})\}
		_{T\tilde{Q}^2}
	\end{align*}
	so that we have
	\begin{align*}
		& \sum_{M = 0}^\infty (-T)^M \{\underbrace{\tilde{X}_s, \ldots,
			\tilde{X}_s}_M\}_{T\tilde{Q}^2} \\
		& \qquad = \sum_{M = 0}^{\infty} (-T)^M \sum_{K=M}^{2M} \sum_{J
			\in \mathcal{P}_{M, K}}
		\{F_{\tilde{\mathscr{M}}_{\mathbb{T}}^{(m)}}(\theta_{J_1}),
		\ldots, F_{\tilde{\mathscr{M}}_{\mathbb{T}}^{(m)}}(\theta_{J_M})\}
		_{T\tilde{Q}^2}
		\\
		& \qquad = \sum_{M = 0}^{\infty} (-T)^M \sum_{K=0}^{M}
		\sum_{J \in \mathcal{P}_{M, K+M}}
		\{F_{\tilde{\mathscr{M}}_{\mathbb{T}}^{(m)}}(\theta_{J_1}),
		\ldots, F_{\tilde{\mathscr{M}}_{\mathbb{T}}^{(m)}}(\theta_{J_M})\}
		_{T\tilde{Q}^2}
		\\
		& \qquad = \sum_{N = 0}^{\infty} \sum_{L=0}^{N}
		(-T)^L \sum_{J \in \mathcal{P}_{L, N}}
		\{F_{\tilde{\mathscr{M}}_{\mathbb{T}}^{(m)}}(\theta_{J_1}),
		\ldots, F_{\tilde{\mathscr{M}}_{\mathbb{T}}^{(m)}}(\theta_{J_L})\}
		_{T\tilde{Q}^2} \\
		& \qquad = \sum_{N=0}^{\infty}
		\Phi_T^{\tilde{\mathscr{M}}_{\mathbb{T}}^{(m)}}
		(\underbrace{\mathcal{A}^s(g),
			\ldots, \mathcal{A}^s(g)}_N)
	\end{align*}
	as desired
\end{proof}
\subsection{Relating the spectral flow to the pairing}\label{Section:
	spectral flow pairing 2}
Since $\omega$ is odd in $\mathrm{Mat}_m(\Omega_{\mathbb{T}})$, we have
\begin{equation*}
	\widetilde{\mathbf{c}_{\mathbb{T}, m}}(\omega) =
	\begin{pmatrix}
		0 & \mathbf{c}_{\mathbb{T}, m}(\omega) \\
		\mathbf{c}_{\mathbb{T}, m}(\omega) & 0
	\end{pmatrix} \in \Lop(\widetilde{\HH^m})
\end{equation*}
where
\begin{equation*}
	\mathbf{c}_{\mathbb{T}, m}(\omega) =
	\begin{pmatrix}
		\mathbf{c}(\omega_{1}^1) &  \dots & \mathbf{c}(\omega_{m}^1) \\
		\vdots & \ddots & \vdots \\
		\mathbf{c}(\omega_{1}^m)  & \dots & \mathbf{c}(\omega_{m}^m)
	\end{pmatrix} \in \Lop(\HH^m).
\end{equation*}
Similarly,
\begin{equation*}
	\widetilde{Q_m} =
	\begin{pmatrix}
		0 & Q_m \\
		Q_m & 0
	\end{pmatrix}
\end{equation*}
on $\widetilde{\HH^m}$ where
\begin{equation*}
	Q_m =
	\begin{pmatrix}
		Q & & 0 \\
		& \ddots & \\
		0  & & Q
	\end{pmatrix}
\end{equation*}
on $\HH^m$.
Therefore,
\begin{align*}
	[\widetilde{Q_m}, \widetilde{\mathbf{c}_{\mathbb{T}, m}}(\omega)] & =
	\widetilde{Q_m}\widetilde{\mathbf{c}_{\mathbb{T}, m}}(\omega) -
	(-1)^{\lvert
		\widetilde{Q_m} \rvert \lvert \widetilde{\mathbf{c}_{\mathbb{T},
				m}}(\omega) \rvert}\widetilde{\mathbf{c}_{\mathbb{T},
			m}}(\omega)
	\widetilde{Q_m} \\
	&=
	\widetilde{Q_m}\widetilde{\mathbf{c}_{\mathbb{T}, m}}(\omega) +
	\widetilde{\mathbf{c}_{\mathbb{T}, m}}(\omega) \widetilde{Q_m} \\
	& =
	\begin{pmatrix}
		[[Q_m, \mathbf{c}_{\mathbb{T}, m}(\omega)]] & 0 \\
		0 & [[Q_m, \mathbf{c}_{\mathbb{T}, m}(\omega)]]
	\end{pmatrix}
\end{align*}
on $\widetilde{\HH^m}$.
Hence, recalling that $U$ is an isomorphism of Fredholm modules, we have
\begin{equation*}
	\tilde{X_s}U = U\big(s[\widetilde{Q_m},
	\widetilde{\mathbf{c}_{\mathbb{T}, m}}(\omega)]
	+ s^2 \widetilde{\mathbf{c}_{\mathbb{T}, m}}(\omega)^2 \big) = U
	\begin{pmatrix}
		X_s & 0 \\
		0 & X_s
	\end{pmatrix}
\end{equation*}
on $\widetilde{\HH^m}$.
Putting everything together, we see that
\begin{align*}
	& \frac 12 \Tr_{\tilde{\HH}^m}\left(\tilde{\gamma}
	\tilde{\mathbf{c}}_{\mathbb{T},
		m}(\omega)e^{-\tau_1\tilde{Q}_{m}^2}\tilde{X}_{s}e^{-(\tau_2
		- \tau_1)\tilde{Q}_{m}^2} \cdot \ldots \cdot e^{-(\tau_{N} -
		\tau_{N-1})\tilde{Q}_{m}^2}\tilde{X}_{s}
	e^{-(1-\tau_{N})\tilde{Q}_{m}^2}
	\right) \\
	& = \frac 12 \Tr_{\widetilde{\HH^m}}\left(U^{-1}\tilde{\gamma}
	\tilde{\mathbf{c}}_{\mathbb{T},
		m}(\omega)e^{-\tau_1\tilde{Q}_{m}^2}\tilde{X}_{s}e^{-(\tau_2
		- \tau_1)\tilde{Q}_{m}^2} \cdot \ldots \cdot e^{-(\tau_{N} -
		\tau_{N-1})\tilde{Q}_{m}^2}\tilde{X}_{s}
	e^{-(1-\tau_{N})\tilde{Q}_{m}^2}
	U\right) \\
	& = \frac 12 \Tr_{\widetilde{\HH^m}}\left(U^{-1}\tilde{\gamma}U
	\widetilde{\mathbf{c}_{\mathbb{T},
			m}}(\omega)e^{-\tau_1\widetilde{Q_m}^2}\mathrm{diag}(X_s) \cdot
	\ldots
	\cdot\mathrm{diag}(X_s)
	e^{-(1-\tau_{N})\widetilde{Q_m}^2}
	\right) \\
	& = \Tr_{\HH^m}\left(\mathbf{c}_{\mathbb{T}, m}(\omega)e^{-\tau_1
		Q_{m}^2}X_{s}e^{-(\tau_2
		- \tau_1) Q_{m}^2} \cdot \ldots \cdot e^{-(\tau_{N} -
		\tau_{N-1}) Q_{m}^2}X_{s}
	e^{-(1-\tau_{N}) Q_{m}^2}
	\right)
\end{align*}
where we used that
\begin{equation*}
	U^{-1} \tilde{\gamma} U =
	\begin{pmatrix}
		0 & \mathbf{1}_m \\
		\mathbf{1}_m & 0
	\end{pmatrix}
\end{equation*}
on $\widetilde{\HH^m}$.
It follows that
\begin{align*}
	& \int_0^1 \Tr_{\HH^m}(\dot{Q}_{g,s} e^{-Q_{g,s}^2}) \, \mathrm{d}s \\
	& \quad = \int_0^1 \Tr_{\HH^m} \left(\sum_{M = 0}^\infty (-1)^M
	\mathbf{c}_{\T, m}(\omega) \{X_s, \ldots, X_s\}_{Q_m^2}\right) \\
	& \quad =
	\int_0^1 \Tr_{\HH^m}\left(\sum_{M = 0}^{\infty}(-1)^M \int_{\Delta_M}
	\mathbf{c}_{\T, m}(\omega)e^{-\tau_1Q_m^2}X_s e^{-(\tau_2
		- \tau_1)Q_m^2} \cdot \ldots \right.\\
	& \qquad \qquad \qquad \qquad \qquad \qquad \left.
	\vphantom{\sum_{M=0}^{\infty}} \ldots
	\cdot
	e^{-(\tau_M -
		\tau_{M-1})Q_m^2}X_s e^{-(1-\tau_M)Q_m^2}\, \mathrm{d} \tau \right)\,
	\mathrm{d}s \\
	& \quad = \int_0^1 \frac 12 \Tr_{\tilde{\HH}^m}\left(\sum_{M =
		0}^{\infty}(-1)^M \int_{\Delta_M}
	\tilde{\gamma} \tilde{\mathbf{c}}_{\mathbb{T},
		m}(\omega)e^{-\tau_1\tilde{Q}_m^2}\tilde{X}_s e^{-(\tau_2
		- \tau_1)\tilde{Q}_m^2} \cdot \ldots \right.\\
	& \qquad \qquad \qquad \qquad \qquad \qquad \left.
	\vphantom{\sum_{M=0}^{\infty}}\ldots \cdot
	e^{-(\tau_M -
		\tau_{M-1})\tilde{Q}_m^2}\tilde{X}_s
	e^{-(1-\tau_M)\tilde{Q}_m^2}\, \mathrm{d} \tau \right)\,
	\mathrm{d}s \\
	& \quad = \int_0^1
	\frac 12 \Tr_{\tilde{\HH}^m}\left(\tilde{\gamma}
	\tilde{\mathbf{c}}_{\mathbb{T},
		m}(\omega) \sum_{M = 0}^{\infty}(-1)^M \{\underbrace{\tilde{X}_s,
		\ldots,
		\tilde{X}_s}_M\}_{\tilde{Q}^2_m} \right)\, \mathrm{d}s \\
	& \quad = \int_0^1
	\frac 12 \Tr_{\tilde{\HH}^m}\left(\tilde{\gamma}
	\tilde{\mathbf{c}}_{\mathbb{T},
		m}(\omega) \sum_{N=0}^{\infty}
	\Phi_1^{\tilde{\mathscr{M}}_{\mathbb{T}}^{(m)}}(\underbrace{\mathcal{A}^s(g),
		\ldots, \mathcal{A}^s(g)}_N) \right)\, \mathrm{d}s \\
	& \quad = \sum_{N=0}^{\infty} \int_0^1 \frac 12
	\Tr_{\tilde{\HH}^m}\Big(\tilde{\gamma}\tilde{\mathbf{c}}_{\mathbb{T},
		m}(\omega)
	\Phi_1^{\tilde{\mathscr{M}}_{\mathbb{T}}^{(m)}}(\mathcal{A}^s(g),
	\ldots, \mathcal{A}^s(g)) \Big)\, \mathrm{d}s
\end{align*}
where we used Fubini's theorem in the last step which is justified by the
estimates in \cref{Fundamental estimate}.
The following statement follows from the proof of \cref{Clifford trace cyclic
sum property}:
\begin{lem}
	For operators $A_1, \ldots, A_{N-1}$ and $B$ satisfying the assumptions
	from \cite[Lemma 4.2]{GüneysuLudewig} and commuting with
	$\tilde{\gamma}$, we have
	\begin{equation*}
		\sum_{k=1}^N \Tr_{\tilde{\HH}^m}(\tilde{\gamma} \{A_k, \ldots,
		A_{N-1},
		B, A_1,
		\ldots, A_{k-1}\}_{\tilde{Q}_m^2}) =
		\Tr_{\tilde{\HH}^m}(\tilde{\gamma} B \{A_1, \ldots,
		A_{N-1}\}_{\tilde{Q}_m^2}).
	\end{equation*}
\end{lem}
By construction, the operators appearing in the brackets of
\begin{equation*}
	\Phi_1^{\tilde{\mathscr{M}}_{\mathbb{T}}^{(m)}}\left(
	\mathcal{A}^s(g)^{\otimes(k-1)} \otimes
	\mathcal{B}^s(g)
	\otimes \mathcal{A}^s(g)^{\otimes(N-k)}
	\right)
\end{equation*}
supercommute with the $\mathcal{C}_1$-action on $\tilde{\HH}^m$ and
therefore, they commute with $\tilde{\gamma}$.
Together with the fact that $F(\mathcal{A}^s(g), \mathcal{B}^s(g)) = 0 =
F(\mathcal{B}^s(g), \mathcal{A}^s(g))$, a similar argument as in
\cref{Ch pullback} shows that
\begin{align*}
	& \Tr_{\tilde{\HH}^m}\left(\tilde{\gamma}
	\sum_{k=1}^{N}\Phi_1^{\tilde{\mathscr{M}}_{\mathbb{T}}^{(m)}}\left(
	\mathcal{A}^s(g)^{\otimes(k-1)} \otimes
	\mathcal{B}^s(g)
	\otimes \mathcal{A}^s(g)^{\otimes(N-k)}
	\right)\right) \\
	& \qquad = \Tr_{\tilde{\HH}^m}\left(\tilde{\gamma} \tilde{\mathbf{c}}_{\T,
		m}\left(\mathcal{B}^s(g)\right)
	\Phi_1^{\tilde{\mathscr{M}}_{\mathbb{T}}^{(m)}}(\mathcal{A}^s(g)^{\otimes(N-1)})
	\right) \\
	& \qquad = \Tr_{\tilde{\HH}^m}\left(\tilde{\gamma}
	\tilde{\mathbf{c}}_{\mathbb{T}, m}(\omega)
	\Phi_1^{\tilde{\mathscr{M}}_{\mathbb{T}}^{(m)}}(\mathcal{A}^s(g)^{\otimes(N-1)})
	\right).
\end{align*}
Substituting this into above calculations, we have
\begin{align*}
	& \int_0^1 \Tr_{\HH^m}(\dot{Q}_{g,s} e^{-Q_{g,s}^2}) \, \mathrm{d}s \\
	& \quad = \sum_{N=0}^{\infty} \int_0^1 \frac 12
	\Tr_{\tilde{\HH}^m}\left(\tilde{\gamma}
	\sum_{k=1}^{N+1}\Phi_1^{\tilde{\mathscr{M}}_{\mathbb{T}}^{(m)}}\left(
	\mathcal{A}^s(g)^{\otimes(k-1)} \otimes
	\mathcal{B}^s(g)
	\otimes \mathcal{A}^s(g)^{\otimes(N+1-k)}
	\right)\right)\, \mathrm{d}s \\
	& \quad = \sum_{N=1}^{\infty} \int_0^1 \frac 12
	\Tr_{\tilde{\HH}^m}\left(\tilde{\gamma}
	\sum_{k=1}^{N}\Phi_1^{\tilde{\mathscr{M}}_{\mathbb{T}}^{(m)}}\left(
	\mathcal{A}^s(g)^{\otimes(k-1)} \otimes
	\mathcal{B}^s(g)
	\otimes \mathcal{A}^s(g)^{\otimes(N-k)}
	\right)\right)\, \mathrm{d}s \\
	& \quad = \sum_{N=0}^{\infty} \int_0^1 \frac 12
	\Tr_{\tilde{\HH}^m}\left(\tilde{\gamma}
	\sum_{k=1}^{N}\Phi_1^{\tilde{\mathscr{M}}_{\mathbb{T}}^{(m)}}\left(
	\mathcal{A}^s(g)^{\otimes(k-1)} \otimes
	\mathcal{B}^s(g)
	\otimes \mathcal{A}^s(g)^{\otimes(N-k)}
	\right)\right)\, \mathrm{d}s
\end{align*}
which is equal to the pairing $\langle
\mathrm{Ch}_{\mathscr{M}_\T}, \mathrm{Ch}(g)\rangle$, see
\cref{Pairing}.
\printbibliography
\end{document}

% $ biblatex auxiliary file $
% $ biblatex bbl format version 3.1 $
% Do not modify the above lines!
%
% This is an auxiliary file used by the 'biblatex' package.
% This file may safely be deleted. It will be recreated by
% biber as required.
%
\begingroup
\makeatletter
\@ifundefined{ver@biblatex.sty}
  {\@latex@error
     {Missing 'biblatex' package}
     {The bibliography requires the 'biblatex' package.}
      \aftergroup\endinput}
  {}
\endgroup


\refsection{0}
  \datalist[entry]{anyt/global//global/global}
    \entry{AtiyahCircular}{inproceedings}{}
      \name{author}{1}{}{%
        {{hash=67ee9912177ae60c68f5fc393dc586be}{%
           family={Atiyah},
           familyi={A\bibinitperiod},
           given={Michael\bibnamedelima F.},
           giveni={M\bibinitperiod\bibinitdelim F\bibinitperiod}}}%
      }
      \name{editor}{2}{}{%
        {{hash=88c5959be0e9e5974a4fc17804884043}{%
           family={Herman},
           familyi={H\bibinitperiod},
           given={Michel\bibnamedelima R.},
           giveni={M\bibinitperiod\bibinitdelim R\bibinitperiod}}}%
        {{hash=e3d9661f68cae7b4fc6e0c7ed5303074}{%
           family={Teissier},
           familyi={T\bibinitperiod},
           given={Bernard},
           giveni={B\bibinitperiod}}}%
      }
      \list{location}{1}{%
        {Palaiseau}%
      }
      \list{publisher}{1}{%
        {As\-té\-risque}%
      }
      \strng{namehash}{67ee9912177ae60c68f5fc393dc586be}
      \strng{fullhash}{67ee9912177ae60c68f5fc393dc586be}
      \strng{bibnamehash}{67ee9912177ae60c68f5fc393dc586be}
      \strng{authorbibnamehash}{67ee9912177ae60c68f5fc393dc586be}
      \strng{authornamehash}{67ee9912177ae60c68f5fc393dc586be}
      \strng{authorfullhash}{67ee9912177ae60c68f5fc393dc586be}
      \strng{editorbibnamehash}{d119942b3d08ed850b9e3174b552fd77}
      \strng{editornamehash}{d119942b3d08ed850b9e3174b552fd77}
      \strng{editorfullhash}{d119942b3d08ed850b9e3174b552fd77}
      \field{labelalpha}{Ati85}
      \field{sortinit}{A}
      \field{sortinithash}{2f401846e2029bad6b3ecc16d50031e2}
      \field{labelnamesource}{author}
      \field{labeltitlesource}{title}
      \field{booktitle}{Colloque en l'honneur de Laurent Schwartz}
      \field{eventyear}{1983}
      \field{number}{131}
      \field{title}{Circular symmetry and stationary-phase approximation}
      \field{venue}{École polytechnique}
      \field{volume}{1}
      \field{year}{1985}
      \field{dateera}{ce}
      \field{eventdateera}{ce}
      \field{pages}{43\bibrangedash 59}
      \range{pages}{17}
      \verb{urlraw}
      \verb http://www.numdam.org/item/AST_1985__131__43_0/
      \endverb
      \verb{url}
      \verb http://www.numdam.org/item/AST_1985__131__43_0/
      \endverb
    \endentry
    \entry{BismutCMP85}{article}{}
      \name{author}{1}{}{%
        {{hash=eace1db2158ab018d2d2bb076c56909b}{%
           family={Bismut},
           familyi={B\bibinitperiod},
           given={Jean-Michel},
           giveni={J\bibinithyphendelim M\bibinitperiod}}}%
      }
      \strng{namehash}{eace1db2158ab018d2d2bb076c56909b}
      \strng{fullhash}{eace1db2158ab018d2d2bb076c56909b}
      \strng{bibnamehash}{eace1db2158ab018d2d2bb076c56909b}
      \strng{authorbibnamehash}{eace1db2158ab018d2d2bb076c56909b}
      \strng{authornamehash}{eace1db2158ab018d2d2bb076c56909b}
      \strng{authorfullhash}{eace1db2158ab018d2d2bb076c56909b}
      \field{labelalpha}{Bis85}
      \field{sortinit}{B}
      \field{sortinithash}{d7095fff47cda75ca2589920aae98399}
      \field{labelnamesource}{author}
      \field{labeltitlesource}{title}
      \field{issue}{2}
      \field{journaltitle}{Communications in Mathematical Physics}
      \field{title}{Index theorem and equivariant cohomology on the loop space}
      \field{volume}{98}
      \field{year}{1985}
      \field{dateera}{ce}
      \field{pages}{213\bibrangedash 237}
      \range{pages}{25}
      \verb{doi}
      \verb https://doi.org/10.1007/BF01220509
      \endverb
    \endentry
    \entry{GüneysuCacciatori}{article}{}
      \name{author}{2}{}{%
        {{hash=aa146ab9abf91df650741771fb0c4e72}{%
           family={Cacciatori},
           familyi={C\bibinitperiod},
           given={Sergio\bibnamedelima L.},
           giveni={S\bibinitperiod\bibinitdelim L\bibinitperiod}}}%
        {{hash=8bc2fa3b59ecd588987a2b2369837761}{%
           family={Güneysu},
           familyi={G\bibinitperiod},
           given={Batu},
           giveni={B\bibinitperiod}}}%
      }
      \strng{namehash}{6c51290fe4800e1bc53421b86e8c2040}
      \strng{fullhash}{6c51290fe4800e1bc53421b86e8c2040}
      \strng{bibnamehash}{6c51290fe4800e1bc53421b86e8c2040}
      \strng{authorbibnamehash}{6c51290fe4800e1bc53421b86e8c2040}
      \strng{authornamehash}{6c51290fe4800e1bc53421b86e8c2040}
      \strng{authorfullhash}{6c51290fe4800e1bc53421b86e8c2040}
      \field{labelalpha}{CG21}
      \field{sortinit}{C}
      \field{sortinithash}{4d103a86280481745c9c897c925753c0}
      \field{labelnamesource}{author}
      \field{labeltitlesource}{title}
      \field{journaltitle}{Journal of Noncommutative Geometry}
      \field{number}{2}
      \field{title}{Odd characteristic classes in entire cyclic homology and equivariant loop space homology}
      \field{volume}{15}
      \field{year}{2021}
      \field{dateera}{ce}
      \field{pages}{615\bibrangedash 642}
      \range{pages}{28}
      \verb{doi}
      \verb https://doi.org/10.4171/jncg/406
      \endverb
    \endentry
    \entry{ConnesCyclicCohomology}{article}{}
      \name{author}{1}{}{%
        {{hash=3745cf9244a5815447aea11906514603}{%
           family={Connes},
           familyi={C\bibinitperiod},
           given={Alain},
           giveni={A\bibinitperiod}}}%
      }
      \strng{namehash}{3745cf9244a5815447aea11906514603}
      \strng{fullhash}{3745cf9244a5815447aea11906514603}
      \strng{bibnamehash}{3745cf9244a5815447aea11906514603}
      \strng{authorbibnamehash}{3745cf9244a5815447aea11906514603}
      \strng{authornamehash}{3745cf9244a5815447aea11906514603}
      \strng{authorfullhash}{3745cf9244a5815447aea11906514603}
      \field{extraname}{1}
      \field{labelalpha}{Con88}
      \field{sortinit}{C}
      \field{sortinithash}{4d103a86280481745c9c897c925753c0}
      \field{labelnamesource}{author}
      \field{labeltitlesource}{title}
      \field{issue}{6}
      \field{journaltitle}{K-Theory}
      \field{month}{11}
      \field{title}{Entire Cyclic Cohomology of Banach Algebras and Characters of $\theta$-Summable Fredholm Modules}
      \field{volume}{1}
      \field{year}{1988}
      \field{dateera}{ce}
      \field{pages}{519\bibrangedash 548}
      \range{pages}{30}
      \verb{doi}
      \verb 10.1007/BF00533785
      \endverb
    \endentry
    \entry{ConnesCompact}{article}{}
      \name{author}{1}{}{%
        {{hash=3745cf9244a5815447aea11906514603}{%
           family={Connes},
           familyi={C\bibinitperiod},
           given={Alain},
           giveni={A\bibinitperiod}}}%
      }
      \strng{namehash}{3745cf9244a5815447aea11906514603}
      \strng{fullhash}{3745cf9244a5815447aea11906514603}
      \strng{bibnamehash}{3745cf9244a5815447aea11906514603}
      \strng{authorbibnamehash}{3745cf9244a5815447aea11906514603}
      \strng{authornamehash}{3745cf9244a5815447aea11906514603}
      \strng{authorfullhash}{3745cf9244a5815447aea11906514603}
      \field{extraname}{2}
      \field{labelalpha}{Con89}
      \field{sortinit}{C}
      \field{sortinithash}{4d103a86280481745c9c897c925753c0}
      \field{labelnamesource}{author}
      \field{labeltitlesource}{title}
      \field{issue}{2}
      \field{journaltitle}{Ergodic Theory and Dynamical Systems}
      \field{month}{6}
      \field{title}{Compact metric spaces, Fredholm modules, and hyperfiniteness}
      \field{volume}{9}
      \field{year}{1989}
      \field{dateera}{ce}
      \field{pages}{207\bibrangedash 220}
      \range{pages}{14}
      \verb{doi}
      \verb https://doi.org/10.1017/S0143385700004934
      \endverb
    \endentry
    \entry{GetzlerOdd}{article}{}
      \name{author}{1}{}{%
        {{hash=96b2c2541889f0975dcf990217ffcda1}{%
           family={Getzler},
           familyi={G\bibinitperiod},
           given={Ezra},
           giveni={E\bibinitperiod}}}%
      }
      \strng{namehash}{96b2c2541889f0975dcf990217ffcda1}
      \strng{fullhash}{96b2c2541889f0975dcf990217ffcda1}
      \strng{bibnamehash}{96b2c2541889f0975dcf990217ffcda1}
      \strng{authorbibnamehash}{96b2c2541889f0975dcf990217ffcda1}
      \strng{authornamehash}{96b2c2541889f0975dcf990217ffcda1}
      \strng{authorfullhash}{96b2c2541889f0975dcf990217ffcda1}
      \field{labelalpha}{Get93}
      \field{sortinit}{G}
      \field{sortinithash}{32d67eca0634bf53703493fb1090a2e8}
      \field{labelnamesource}{author}
      \field{labeltitlesource}{title}
      \field{issn}{0040-9383}
      \field{journaltitle}{Topology}
      \field{number}{3}
      \field{title}{The odd Chern character in cyclic homology and spectral flow}
      \field{volume}{32}
      \field{year}{1993}
      \field{pages}{489\bibrangedash 507}
      \range{pages}{19}
      \verb{doi}
      \verb https://doi.org/10.1016/0040-9383(93)90002-D
      \endverb
    \endentry
    \entry{GetzlerJonesPetrack}{article}{}
      \name{author}{3}{}{%
        {{hash=96b2c2541889f0975dcf990217ffcda1}{%
           family={Getzler},
           familyi={G\bibinitperiod},
           given={Ezra},
           giveni={E\bibinitperiod}}}%
        {{hash=fdbc450b184f7e09d25d26162515d42c}{%
           family={Jones},
           familyi={J\bibinitperiod},
           given={John\bibnamedelimb D.\bibnamedelimi S.},
           giveni={J\bibinitperiod\bibinitdelim D\bibinitperiod\bibinitdelim S\bibinitperiod}}}%
        {{hash=84663be93d081a473d7459ab66c541dd}{%
           family={Petrack},
           familyi={P\bibinitperiod},
           given={Scott},
           giveni={S\bibinitperiod}}}%
      }
      \strng{namehash}{586685fc2b0dc06950da11d9fa0ea9a4}
      \strng{fullhash}{586685fc2b0dc06950da11d9fa0ea9a4}
      \strng{bibnamehash}{586685fc2b0dc06950da11d9fa0ea9a4}
      \strng{authorbibnamehash}{586685fc2b0dc06950da11d9fa0ea9a4}
      \strng{authornamehash}{586685fc2b0dc06950da11d9fa0ea9a4}
      \strng{authorfullhash}{586685fc2b0dc06950da11d9fa0ea9a4}
      \field{labelalpha}{GJP91}
      \field{sortinit}{G}
      \field{sortinithash}{32d67eca0634bf53703493fb1090a2e8}
      \field{labelnamesource}{author}
      \field{labeltitlesource}{title}
      \field{issue}{3}
      \field{journaltitle}{Topology}
      \field{month}{11}
      \field{title}{Differential forms on loop spaces and the cyclic bar complex}
      \field{volume}{30}
      \field{year}{1991}
      \field{dateera}{ce}
      \field{pages}{339\bibrangedash 371}
      \range{pages}{33}
      \verb{doi}
      \verb https://doi.org/10.1016/0040-9383(91)90019-Z
      \endverb
    \endentry
    \entry{GüneysuLudewig}{article}{}
      \name{author}{2}{}{%
        {{hash=8bc2fa3b59ecd588987a2b2369837761}{%
           family={Güneysu},
           familyi={G\bibinitperiod},
           given={Batu},
           giveni={B\bibinitperiod}}}%
        {{hash=1f9f7f3b43c82f3796cdbd6f1e61849f}{%
           family={Ludewig},
           familyi={L\bibinitperiod},
           given={Matthias},
           giveni={M\bibinitperiod}}}%
      }
      \strng{namehash}{817d4462ca7b60312fb9edd68226890a}
      \strng{fullhash}{817d4462ca7b60312fb9edd68226890a}
      \strng{bibnamehash}{817d4462ca7b60312fb9edd68226890a}
      \strng{authorbibnamehash}{817d4462ca7b60312fb9edd68226890a}
      \strng{authornamehash}{817d4462ca7b60312fb9edd68226890a}
      \strng{authorfullhash}{817d4462ca7b60312fb9edd68226890a}
      \field{labelalpha}{GL22}
      \field{sortinit}{G}
      \field{sortinithash}{32d67eca0634bf53703493fb1090a2e8}
      \field{labelnamesource}{author}
      \field{labeltitlesource}{title}
      \field{issn}{0001-8708}
      \field{journaltitle}{Advances in Mathematics}
      \field{number}{108143}
      \field{title}{The Chern character of $\vartheta$-summable Fredholm modules over dg algebras and localization on loop space}
      \field{volume}{395}
      \field{year}{2022}
      \verb{doi}
      \verb https://doi.org/10.1016/j.aim.2021.108143
      \endverb
    \endentry
    \entry{GetzlerSzenes}{article}{}
      \name{author}{2}{}{%
        {{hash=96b2c2541889f0975dcf990217ffcda1}{%
           family={Getzler},
           familyi={G\bibinitperiod},
           given={Ezra},
           giveni={E\bibinitperiod}}}%
        {{hash=cdea49209331300fcef8df31ffc06da4}{%
           family={Szenes},
           familyi={S\bibinitperiod},
           given={András},
           giveni={A\bibinitperiod}}}%
      }
      \strng{namehash}{a471202534002d704bb28a797a3f31a2}
      \strng{fullhash}{a471202534002d704bb28a797a3f31a2}
      \strng{bibnamehash}{a471202534002d704bb28a797a3f31a2}
      \strng{authorbibnamehash}{a471202534002d704bb28a797a3f31a2}
      \strng{authornamehash}{a471202534002d704bb28a797a3f31a2}
      \strng{authorfullhash}{a471202534002d704bb28a797a3f31a2}
      \field{labelalpha}{GS89}
      \field{sortinit}{G}
      \field{sortinithash}{32d67eca0634bf53703493fb1090a2e8}
      \field{labelnamesource}{author}
      \field{labeltitlesource}{title}
      \field{issue}{2}
      \field{journaltitle}{Journal of Functional Analysis}
      \field{month}{6}
      \field{title}{On the Chern character of a theta-summable Fredholm module}
      \field{volume}{84}
      \field{year}{1989}
      \field{dateera}{ce}
      \field{pages}{343\bibrangedash 357}
      \range{pages}{15}
      \verb{doi}
      \verb https://doi.org/10.1016/0022-1236(89)90102-X
      \endverb
    \endentry
    \entry{JLO}{article}{}
      \name{author}{3}{}{%
        {{hash=d685c6cc3774c2c910a76dcd7da5bf36}{%
           family={Jaffe},
           familyi={J\bibinitperiod},
           given={Arthur},
           giveni={A\bibinitperiod}}}%
        {{hash=106662116187b2f35877a0d17b073696}{%
           family={Lesniewski},
           familyi={L\bibinitperiod},
           given={Andrzej},
           giveni={A\bibinitperiod}}}%
        {{hash=4a9a5bdcc30010dc439c8dc664dd6302}{%
           family={Osterwalder},
           familyi={O\bibinitperiod},
           given={Konrad},
           giveni={K\bibinitperiod}}}%
      }
      \strng{namehash}{0f1e8a233905d2ec5b5eedac93d515e9}
      \strng{fullhash}{0f1e8a233905d2ec5b5eedac93d515e9}
      \strng{bibnamehash}{0f1e8a233905d2ec5b5eedac93d515e9}
      \strng{authorbibnamehash}{0f1e8a233905d2ec5b5eedac93d515e9}
      \strng{authornamehash}{0f1e8a233905d2ec5b5eedac93d515e9}
      \strng{authorfullhash}{0f1e8a233905d2ec5b5eedac93d515e9}
      \field{labelalpha}{JLO88}
      \field{sortinit}{J}
      \field{sortinithash}{b2f54a9081ace9966a7cb9413811edb4}
      \field{labelnamesource}{author}
      \field{labeltitlesource}{title}
      \field{journaltitle}{Communications in Mathematical Physics}
      \field{number}{1}
      \field{subtitle}{I. The Chern character}
      \field{title}{Quantum K-Theory}
      \field{volume}{118}
      \field{year}{1988}
      \field{dateera}{ce}
      \field{pages}{1\bibrangedash 14}
      \range{pages}{14}
      \verb{doi}
      \verb https://doi.org/10.1007/BF01218474
      \endverb
    \endentry
    \entry{Kato}{book}{}
      \name{author}{1}{}{%
        {{hash=de2df565ad24725681a931fe6bd95efd}{%
           family={Kato},
           familyi={K\bibinitperiod},
           given={Tosio},
           giveni={T\bibinitperiod}}}%
      }
      \list{publisher}{1}{%
        {Springer-Verlag}%
      }
      \strng{namehash}{de2df565ad24725681a931fe6bd95efd}
      \strng{fullhash}{de2df565ad24725681a931fe6bd95efd}
      \strng{bibnamehash}{de2df565ad24725681a931fe6bd95efd}
      \strng{authorbibnamehash}{de2df565ad24725681a931fe6bd95efd}
      \strng{authornamehash}{de2df565ad24725681a931fe6bd95efd}
      \strng{authorfullhash}{de2df565ad24725681a931fe6bd95efd}
      \field{labelalpha}{Kat95}
      \field{sortinit}{K}
      \field{sortinithash}{c02bf6bff1c488450c352b40f5d853ab}
      \field{labelnamesource}{author}
      \field{labeltitlesource}{title}
      \field{edition}{Corrected Printing of the Second Edition}
      \field{number}{132}
      \field{series}{Classics in Mathematics}
      \field{title}{Perturbation Theory for Linear Operators}
      \field{year}{1995}
      \field{dateera}{ce}
      \verb{doi}
      \verb https://doi.org/10.1007/978-3-642-66282-9
      \endverb
    \endentry
  \enddatalist
\endrefsection
\endinput

